% ============================================================
%  Separation Technologies and Techniques — Experiment Booklet (one section per experiment)
%  Year 7 (Year 8 extension)  |  Victorian Curriculum F–10 Version 2.0
%  Compile with: xelatex three times, then python3 fill.py Separation_Experiment_Booklet_Yr7
% ============================================================
\documentclass[11pt,a4paper]{article}
\newcommand{\unitname}{Separation Technologies and Techniques}\newcommand{\unitcaps}{SEPARATION}

% ── Fonts (compile with LuaLaTeX) ────────────────────────────────────────────
\usepackage{fontspec}
% Bundled fonts (in ./fonts/) so the document compiles anywhere with XeLaTeX/LuaLaTeX
\setmainfont{Lato}[                   % body — humanist sans (stands in for deck's Inter)
  Path=./fonts/, Extension=.ttf, Ligatures=TeX,
  UprightFont=Lato-Regular, ItalicFont=Lato-Italic,
  BoldFont=Lato-Bold, BoldItalicFont=Lato-BoldItalic ]
\newfontfamily\headingfont{Poppins}[  % display/headings — matches the slide decks
  Path=./fonts/, Extension=.ttf, Ligatures=TeX,
  UprightFont=Poppins-Regular, ItalicFont=Poppins-Italic,
  BoldFont=Poppins-Bold, BoldItalicFont=Poppins-BoldItalic ]
\setlength{\headheight}{22pt}

\usepackage[protrusion=true,expansion=false]{microtype}
\usepackage[a4paper, top=2cm, bottom=2cm, left=2cm, right=2cm]{geometry}
\usepackage{xcolor}
\usepackage{tikz}
\usetikzlibrary{arrows.meta, shapes.geometric, calc, positioning, decorations.pathreplacing}
\usepackage{tcolorbox}
\tcbuselibrary{skins, breakable, raster}
\usepackage{fancyhdr}
\usepackage{titlesec}
\usepackage{enumitem}
\usepackage{booktabs}
\usepackage{tabularx}
\usepackage{array}
\usepackage{colortbl}
\usepackage{multicol}
\usepackage{amsmath}
\usepackage{needspace}
% TikZ diagram macros are embedded below (just before \begin{document}) for easy editing.
\usepackage{longtable}
\usepackage{lastpage}
\usepackage{calc}
\usepackage{mdframed}
\usepackage{setspace}
\usepackage{ragged2e}
% PDF bookmarks: one per lesson (made in \sectiondiv) plus the cover and
% contents, and links from the contents table to each lesson.
% implicit=false: hyperref adds no anchors of its own (an anchor on each
% list item nudged lists down 1.5 mm, which would undo the page fill).
% hidelinks keeps the links invisible, so the printed pages are unchanged.
\usepackage[implicit=false, hidelinks,
  pdftitle={Separation Technologies and Techniques — Experiment Booklet},
  pdfsubject={Years 7–8 Science, Victorian Curriculum F–10 v2.0}]{hyperref}
\usepackage{bookmark} % writes the outline itself, which implicit=false switches off in hyperref
% hyperref 7.01 with implicit=false leaves \Hy@PageAnchorSlide undefined, which
% throws at \end{document} and stops fill.py (it builds with -halt-on-error).
\makeatletter\providecommand\Hy@PageAnchorSlide{}\makeatother

% ── Colour Palette ───────────────────────────────────────────────────────────
% Remapped to the slide-deck palette (names kept so body + diagrams cascade)
\definecolor{NavyBlue}   {HTML}{16305C}
\definecolor{MidBlue}    {HTML}{2E6FE0}
\definecolor{LightBlue}  {HTML}{E8F1FF}
\definecolor{TealGreen}  {HTML}{0F6E63}
\definecolor{LightGreen} {HTML}{E0F7F3}
\definecolor{GoldDark}   {HTML}{C77F12}
\definecolor{GoldLight}  {HTML}{FFF3DC}
\definecolor{PurpleDark} {HTML}{5B3FC0}
\definecolor{PurpleLight}{HTML}{EFE9FF}
\definecolor{OrangeDark} {HTML}{C9483D}
\definecolor{OrangeLight}{HTML}{FFE9E6}
\definecolor{MidGrey}    {HTML}{AEBED4}
\definecolor{LightGrey}  {HTML}{F2F5FA}
\definecolor{DarkGrey}   {HTML}{3A4759}
\definecolor{TableHead}  {HTML}{16305C}
\definecolor{RowAlt}     {HTML}{EEF4FF}
% Bright accent tints (deck) for bars, rules and chips
\definecolor{Blue}   {HTML}{2E6FE0}
\definecolor{Teal}   {HTML}{18A999}
\definecolor{Amber}  {HTML}{F5A623}
\definecolor{Coral}  {HTML}{FF6B5E}
\definecolor{Violet} {HTML}{7C5CE0}
\definecolor{Sky}    {HTML}{4DA3FF}
\definecolor{Ink}    {HTML}{16223A}
\definecolor{Line}   {HTML}{DDE5F0}
\definecolor{Paper}  {HTML}{FBFCFE}

% ── Page Style ───────────────────────────────────────────────────────────────
\pagestyle{fancy}
\fancyhf{}
\renewcommand{\headrulewidth}{0.4pt}
\renewcommand{\headrule}{\hbox to\headwidth{\color{NavyBlue}\leaders\hrule height \headrulewidth\hfill}}
\fancyhead[L]{\headingfont\small\color{NavyBlue}\unitname{} \textbar{} Experiment Booklet}
\fancyhead[R]{\headingfont\small\color{DarkGrey}Year 7 \textbar{} Victorian Curriculum F--10 v2.0}
\fancyfoot[C]{\headingfont\small\color{DarkGrey}\thepage\ \textbar{} \pageref{LastPage}}
\fancypagestyle{plain}{%
  \fancyhf{}
  \fancyfoot[C]{\headingfont\small\color{DarkGrey}\thepage\ \textbar{} \pageref{LastPage}}
  \renewcommand{\headrulewidth}{0pt}
}

% ── Section Formatting ───────────────────────────────────────────────────────
\titleformat{\section}
  {\headingfont\color{NavyBlue}\Large\bfseries}
  {}{0em}{}
  [{\color{Teal}\rule[0.6ex]{2.2em}{3pt}}{\color{Line}\rule[0.6ex]{\dimexpr\linewidth-2.2em\relax}{1pt}}]
\titlespacing{\section}{0pt}{18pt}{8pt}

\titleformat{\subsection}
  {\headingfont\color{MidBlue}\large\bfseries}
  {\thesubsection}{0.6em}{}
\titlespacing{\subsection}{0pt}{12pt}{4pt}

\titleformat{\subsubsection}
  {\headingfont\color{NavyBlue}\normalsize\bfseries}
  {}{0em}{}
\titlespacing{\subsubsection}{0pt}{8pt}{2pt}

% Deck-style eyebrow / kicker (optional use)
\newcommand{\kicker}[2][Teal]{%
  {\headingfont\bfseries\footnotesize\color{#1}%
   \textcolor{#1}{\rule[0.18ex]{1.6em}{2.2pt}}\hspace{0.5em}\MakeUppercase{#2}\par}%
  \vspace{2pt}}

% ── Custom tcolorbox Environments ────────────────────────────────────────────

% Section divider banner
\newtcolorbox{sectionbanner}[2]{
  enhanced, breakable=false,
  colback=NavyBlue, colframe=NavyBlue,
  fontupper=\color{white}\bfseries\large,
  left=4mm, right=4mm, top=3mm, bottom=3mm,
  title={\color{white}\Huge\bfseries #1},
  attach boxed title to top left={yshift=-2mm, xshift=4mm},
  boxed title style={colback=NavyBlue, colframe=NavyBlue},
  after upper={\par\smallskip\normalfont\normalsize\color{white!80!NavyBlue}\textit{#2}}
}

% Key info box (blue)
\newtcolorbox{infobox}[1]{
  breakable, arc=2.5mm, boxrule=0.5pt,
  colback=LightBlue, colframe=Blue!40,
  fonttitle=\headingfont\bfseries, coltitle=MidBlue,
  title={#1},
  left=4mm, right=3.5mm, top=2.5mm, bottom=2.5mm
}

% Scaffold box (yellow)
\newtcolorbox{scaffold}[1]{
  breakable, arc=2.5mm, boxrule=0.5pt,
  colback=GoldLight, colframe=Amber!40,
  fonttitle=\headingfont\bfseries, coltitle=GoldDark,
  title={#1},
  left=4mm, right=3.5mm, top=2.5mm, bottom=2.5mm
}

% Worked example box
\newtcolorbox{workedexbox}[1]{
  breakable, arc=2.5mm, boxrule=0.5pt,
  colback=white, colframe=Amber!50,
  fonttitle=\headingfont\bfseries, coltitle=GoldDark,
  title={#1},
  left=4mm, right=3.5mm, top=2.5mm, bottom=2.5mm
}
% Worked examples: \begin{workedex}{title} -> "Worked Example N.M --- title",
% auto-numbered per section and labelled we:N.M. Never type the number.
\newcounter{wenum}[section]
\renewcommand{\thewenum}{\thesection.\arabic{wenum}}
\newenvironment{workedex}[1]
  {\refstepcounter{wenum}\label{we:\thewenum}\begin{workedexbox}{Worked Example \thewenum\ --- #1}}
  {\end{workedexbox}}

% We do: straight after each worked example, a near-identical question the
% class works through together, with its own writing space. Numbered to
% match the worked example above it and labelled wd:N.M.
%   \begin{wedo}{title} question ... \answerbox{h} \end{wedo}
\newtcolorbox{wedobox}[1]{
  enhanced, breakable=false, arc=2.5mm, boxrule=0.5pt,
  colback=white, colframe=TealGreen!55,
  fonttitle=\headingfont\bfseries, coltitle=TealGreen!75!black,
  title={#1},
  left=4mm, right=3.5mm, top=2.5mm, bottom=2.5mm
}
\newenvironment{wedo}[1]
  {\par\needspace{9\baselineskip}\begin{wedobox}{We do \thewenum\ --- #1}\label{wd:\thewenum}\ignorespaces}
  {\end{wedobox}}

% Year 7 box (blue)
\newtcolorbox{yr7box}{
  breakable, arc=2.5mm, boxrule=0.5pt,
  colback=LightBlue, colframe=Blue!40,
  left=4mm, right=3.5mm, top=2.5mm, bottom=2.5mm
}

% Year 8 extension box (violet)
\newtcolorbox{yr8box}{
  breakable, arc=2.5mm, boxrule=0.5pt,
  colback=PurpleLight, colframe=Violet!40,
  left=4mm, right=3.5mm, top=2.5mm, bottom=2.5mm
}

% Answer lines box
\newcommand{\answerbox}[1]{%
  \begin{tcolorbox}[
    enhanced, breakable, arc=2mm,
    colback=Paper, colframe=Line,
    left=2mm, right=2mm, top=1mm, bottom=1mm,
    boxrule=1pt,
    height=#1
  ]
  \end{tcolorbox}%
}

% Experiment header
\newtcolorbox{expheader}[3]{
  breakable=false, arc=2.5mm,
  colback=RowAlt, colframe=Amber!50, boxrule=0.5pt,
  fonttitle=\headingfont\bfseries\large, coltitle=NavyBlue,
  title={Experiment #1: #2},
  left=4mm, right=3.5mm, top=2.5mm, bottom=2.5mm,
  after upper={\par\smallskip
    {\headingfont\bfseries\color{GoldDark}Aim:}\ \textit{#3}},
}

% Reading comprehension passage box (green)
\newtcolorbox{readingbox}[1]{
  breakable, arc=2.5mm, boxrule=0.5pt,
  colback=white, colframe=Teal!50,
  fonttitle=\headingfont\bfseries, coltitle=TealGreen,
  title={#1},
  left=4mm, right=3.5mm, top=2.5mm, bottom=2.5mm,
  fontupper=\normalsize
}

% Skill-builder box (for beginner scaffolding)
\newtcolorbox{skillbox}[1]{
  breakable, arc=2.5mm, boxrule=0.5pt,
  colback=RowAlt, colframe=Blue!40,
  fonttitle=\headingfont\bfseries, coltitle=MidBlue,
  title={#1},
  left=4mm, right=3.5mm, top=2.5mm, bottom=2.5mm
}

% ── Custom Commands ───────────────────────────────────────────────────────────

% Question: \question{marks}{text}  (marks can be empty)
% Numbered automatically as Q<section>.<n>, reset every section.
% Do NOT write numbers by hand.
\newcounter{qnum}[section]
\renewcommand{\theqnum}{\thesection.\arabic{qnum}}
\newcommand{\question}[2]{%
  \refstepcounter{qnum}%
  \needspace{3\baselineskip}%
  \medskip
  \noindent\label{q:\theqnum}\textbf{Q\theqnum}\ifx&#1&\else\ \textcolor{DarkGrey}{(#1 mark\ifnum#1=1\else s\fi)}\fi\quad #2%
  \medskip
}

% Inline marks note
\newcommand{\markscount}[1]{\hfill\textcolor{DarkGrey}{\textit{(#1 mark\ifnum#1=1\else s\fi)}}}

% ── Practice set counter and question macros ─────────────────────────────────
\newcounter{pqnum}
\newcommand{\practiceset}[1]{%
  \needspace{6\baselineskip}
  \setcounter{pqnum}{0}%
  \vspace{6pt}
  {\headingfont\bfseries\large\color{GoldDark} Practice Set --- #1}\\[-2pt]
  {\color{Amber}\rule{\linewidth}{1pt}}\\[2pt]
}
% \pq{question text}{working box height}
\newcommand{\pq}[2]{%
  \stepcounter{pqnum}%
  \needspace{4\baselineskip}
  \vspace{5pt}
  \noindent\textbf{\thepqnum.}\quad #1\par
  \vspace{2pt}
  \answerbox{#2}
}
% Two short questions side by side, each with its own box
\newcommand{\pqpair}[4]{%
  \needspace{4\baselineskip}
  \vspace{5pt}
  \noindent
  \begin{minipage}[t]{0.48\linewidth}
    \stepcounter{pqnum}\textbf{\thepqnum.}\quad #1\par\vspace{2pt}
    \begin{tcolorbox}[enhanced,arc=2mm,colback=Paper,colframe=Line,left=2mm,right=2mm,
      top=1mm,bottom=1mm,boxrule=1pt,height=#2]\end{tcolorbox}
  \end{minipage}\hfill
  \begin{minipage}[t]{0.48\linewidth}
    \stepcounter{pqnum}\textbf{\thepqnum.}\quad #3\par\vspace{2pt}
    \begin{tcolorbox}[enhanced,arc=2mm,colback=Paper,colframe=Line,left=2mm,right=2mm,
      top=1mm,bottom=1mm,boxrule=1pt,height=#4]\end{tcolorbox}
  \end{minipage}
}

% ── Comprehension question macro (compact, lined) ────────────────────────────
\newcounter{rcnum}
\newcommand{\rcreset}{\setcounter{rcnum}{0}}
\newcommand{\rcq}[2]{%
  \stepcounter{rcnum}%
  \needspace{3\baselineskip}
  \vspace{4pt}
  \noindent\textbf{C\thercnum.}\quad #1\par
  \vspace{2pt}
  \linedbox{#2}
}

% Section divider (big numbered banner): \sectiondiv{title}{summary}
%
% This IS the section: it steps the real section counter, so subsections,
% tasks, questions and figures inside it are numbered N.M automatically.
% Put \label{sec:N} straight after it. Never hand-type a section number.
\newcommand{\sectiondiv}[2]{%
  \needspace{5\baselineskip}%
  \refstepcounter{section}%
  % PDF bookmark for the lesson; its anchor, lesson.N.1, is what the contents table links to
  \pdfbookmark[1]{Lesson \thesection{} — #1}{lesson.\thesection}%
  \begin{tcolorbox}[
    enhanced, breakable=false, arc=3.5mm,
    colback=NavyBlue, colframe=NavyBlue, boxrule=0pt,
    sidebyside, sidebyside align=center,
    lefthand width=2.4cm,
    left=2mm, right=4mm, top=4mm, bottom=4mm
  ]
    \begin{center}
      {\headingfont\bfseries\color{white!35!NavyBlue}\fontsize{50}{50}\selectfont \thesection}
    \end{center}
    \tcblower
    {\headingfont\bfseries\color{white}\LARGE #1}\\[3pt]
    {\itshape\color{white!75!NavyBlue}\small #2}
  \end{tcolorbox}%
}

% Tasks: \task, \task[a] (part a of a NEW task), \taskpart{b} (next part
% of the same task), \exttask (Year 8 "Extension" — same counter).
% Numbered Task <section>.<n>, reset every section. Do NOT type numbers.
% Every task, question and figure also gets an automatic label
% (task:1.7b, q:2.4, fig:1.2, fig:E2.1) so page refs come from the .aux.
\newcounter{tasknum}[section]
\renewcommand{\thetasknum}{\thesection.\arabic{tasknum}}
\NewDocumentCommand{\task}{o}{\refstepcounter{tasknum}\label{task:\thetasknum\IfValueT{#1}{#1}}\textbf{Task \thetasknum\IfValueT{#1}{#1}}}
\newcommand{\taskpart}[1]{\label{task:\thetasknum#1}\textbf{Task \thetasknum#1}}
\newcommand{\exttask}{\refstepcounter{tasknum}\label{task:\thetasknum}\textbf{Extension \thetasknum}}

% Figure captions: \figcap[fig:label]{caption} -> "Figure N.M --- caption"
% \expfigcap{caption} -> "Figure EN.M" (experiment pages). Both auto-numbered.
\newcounter{fignum}[section]
\renewcommand{\thefignum}{\thesection.\arabic{fignum}}
\newcounter{expfignum}[section]
\renewcommand{\theexpfignum}{E\thesection.\arabic{expfignum}}
\newcommand{\figcap}[2][]{%
  \begin{center}
    \refstepcounter{fignum}\label{fig:\thefignum}\if\relax\detokenize{#1}\relax\else\label{#1}\fi
    \small\itshape\color{DarkGrey} Figure \thefignum\ --- #2
  \end{center}%
}
\newcommand{\expfigcap}[1]{%
  \begin{center}
    \refstepcounter{expfignum}\label{fig:\theexpfignum}%
    \small\itshape\color{DarkGrey} Figure \theexpfignum\ --- #1
  \end{center}%
}

% Diagram include
\newcommand{\diagramfig}[3]{%  {filename}{width fraction}{caption}
  \begin{center}
    \includegraphics[width=#2\linewidth]{#1}\\[4pt]
    \small\itshape\color{DarkGrey} #3
  \end{center}%
}

% Table column types
\newcolumntype{L}[1]{>{\RaggedRight\arraybackslash}p{#1}}
\newcolumntype{C}[1]{>{\Centering\arraybackslash}p{#1}}

% Table header colour
\newcommand{\thead}[1]{\cellcolor{TableHead}\color{white}\bfseries #1}
\newcommand{\talt}{\rowcolor{RowAlt}}
\newcommand{\tmid}{\rowcolor{MidGrey}}

% Inline year tags
\newcommand{\tagyr}[1]{%
  \tikz[baseline=-0.5ex]\node[
    rounded corners=2pt,
    fill=#1!20, draw=#1, text=#1,
    inner xsep=4pt, inner ysep=1pt,
    font=\scriptsize\bfseries
  ]{Yr #1 only};}

% Draw answer lines
\newcommand{\anslines}[1]{%
  \vspace{2pt}%
  \foreach \n in {1,...,#1}{%
    \par\noindent\rule{\linewidth}{0.3pt}\vspace{3pt}%
  }%
}

% ── Document ─────────────────────────────────────────────────────────────────
% ── Optional-lesson note, shown under the lesson banner ──────────────────────
\newcommand{\optionalnote}[1]{%
  % the title is a bold lead-in on the first line rather than a title bar,
  % which keeps Lesson 11's reading on the lesson's first page
  \begin{tcolorbox}[enhanced, breakable=false, arc=2.5mm, boxrule=0.5pt,
    colback=PurpleLight, colframe=Violet!50,
    left=4mm, right=3.5mm, top=2mm, bottom=2mm]
  {\headingfont\bfseries\color{PurpleDark}Optional lesson \textbullet{} Year 8 extension.}\enspace #1
  \end{tcolorbox}\vspace{2pt}}


% ══════════════════════════════════════════════════════════════════════════════
%  LESSON STRUCTURE (lesson booklet, Oct 2026)
%  Every lesson runs in the same four parts, in this order:
%    \lessonpart{A}{Reading and comprehension}{Teal}   \label{read:N}
%    \lessonpart{B}{Theory and examples}{Blue}         \label{theory:N}
%    \lessonpart{C}{Practice}{Amber}                    \label{prac:N}
%    \lessonpart{D}{Above level}{Violet}                \label{ext:N}
%  Above-level / higher-order questions live ONLY in Part D, at the end.
% ══════════════════════════════════════════════════════════════════════════════
\newcommand{\lessonpart}[3]{%
  \par\if#1A\needspace{10\baselineskip}\vspace{10pt}\else
    \if#1C\notesfill\fi\clearpage\fi
  \begin{tcolorbox}[enhanced, breakable=false, arc=2mm, boxrule=0pt,
    colback=#3!14, borderline west={4pt}{0pt}{#3},
    left=4mm, right=3mm, top=1.6mm, bottom=1.6mm]
  {\headingfont\bfseries\color{#3!60!black}\large Part #1 \textbullet{} #2}
  \end{tcolorbox}\vspace{2pt}}

% Key vocabulary at the start of each lesson: students copy each definition
% from the vocabulary slide in the lesson's starter. First row = red scaffold.
\newcommand{\vocabhead}{%
  \par\vspace{6pt}%
  \begin{tcolorbox}[enhanced, breakable=false, arc=2mm, boxrule=0pt,
    colback=NavyBlue!10, borderline west={4pt}{0pt}{NavyBlue},
    left=4mm, right=3mm, top=1.6mm, bottom=1.6mm]
  {\headingfont\bfseries\color{NavyBlue}\large Key vocabulary}
  \end{tcolorbox}\vspace{2pt}%
  \noindent Copy the definition of each word from the slide. The first one has been done for you.\par}

% Curriculum codes under a lesson banner
\newcommand{\lessoncodes}[1]{%
  {\small\headingfont\color{DarkGrey}Victorian Curriculum 2.0:\ \textbf{#1}}\par\vspace{4pt}}

% Two short tasks side by side, each with its own working box
% \taskpair{text}{box height}{text}{box height}
% The two boxes line up: the shorter question is padded by the difference in
% text height (measured in a scratch box, task counter restored, no label).
\newsavebox\tpbox
\newlength\tpA \newlength\tpB
\newcount\tpsaved
\newcommand\tpmeasure[2]{%
  \tpsaved=\value{tasknum}%
  \sbox\tpbox{\vbox{\hsize=0.485\linewidth \linewidth=\hsize \parskip=0pt
    \renewcommand\label[1]{}\task #1\par}}%
  \setlength{#2}{\dimexpr\ht\tpbox+\dp\tpbox\relax}%
  \setcounter{tasknum}{\tpsaved}}
\newcommand{\taskpair}[4]{%
  \par\needspace{9\baselineskip}\vspace{4pt}%
  \tpmeasure{#1}{\tpA}\tpmeasure{#3}{\tpB}%
  \noindent
  \edef\abgrp{\number\numexpr\value{abox}+1\relax}%
  \begin{minipage}[t]{0.485\linewidth}
    \task #1\par\ifdim\tpB>\tpA\vspace{\dimexpr\tpB-\tpA\relax}\fi\vspace{2pt}\answerbox{#2}
  \end{minipage}\hfill
  \begin{minipage}[t]{0.485\linewidth}
    \task #3\par\ifdim\tpA>\tpB\vspace{\dimexpr\tpA-\tpB\relax}\fi\vspace{2pt}\answerbox{#4}
  \end{minipage}\par\gdef\abgrp{0}\vspace{4pt}}
% Two parts of one question side by side, with their answer boxes level:
% \ptpair{a}{part (a) text}{box height}{b}{part (b) text}{box height}.
% The shorter part's text gets space below it to match the longer one.
\newcommand\ptmeasure[3]{%
  \sbox\tpbox{\vbox{\hsize=0.485\linewidth \linewidth=\hsize \parskip=0pt \pt{#1}#2\par}}%
  \setlength{#3}{\dimexpr\ht\tpbox+\dp\tpbox\relax}}
\newcommand{\ptpair}[6]{%
  \ptmeasure{#1}{#2}{\tpA}\ptmeasure{#4}{#5}{\tpB}%
  \par\noindent
  \edef\abgrp{\number\numexpr\value{abox}+1\relax}%
  \begin{minipage}[t]{0.485\linewidth}
    \pt{#1}#2\par\ifdim\tpB>\tpA\vspace{\dimexpr\tpB-\tpA\relax}\fi\answerbox{#3}
  \end{minipage}\hfill
  \begin{minipage}[t]{0.485\linewidth}
    \pt{#4}#5\par\ifdim\tpA>\tpB\vspace{\dimexpr\tpA-\tpB\relax}\fi\answerbox{#6}
  \end{minipage}\par\gdef\abgrp{0}}

% Answer boxes never split across a page (a half box on its own page is useless)
% ── Page filling ────────────────────────────────────────────────────────────
% Every writing box gets an id. Its page is logged to \jobname.abp; fill.py
% measures the empty space left at the foot of each page and shares it out
% among that page's boxes via \jobname.abx (\abset{id}{extra}). Side-by-side
% boxes share a group so they grow together.
\newcounter{abox}
\newwrite\abpos \immediate\openout\abpos=\jobname.abp
\newcommand{\abset}[2]{\expandafter\gdef\csname abx@#1\endcsname{#2}}
\InputIfFileExists{\jobname.abx}{}{}
\def\abgrp{0}
\newcommand{\abnext}{\stepcounter{abox}%
  \edef\abcur{\ifcsname abx@\theabox\endcsname\csname abx@\theabox\endcsname\else 0pt\fi}}
\newcommand{\ablog}{\edef\abtmp{\noexpand\write\abpos{\theabox\space\abgrp\space\noexpand\thepage}}\abtmp}
\newcommand{\abh}[1]{\dimexpr#1+\abcur\relax}
\renewcommand{\answerbox}[1]{%
  \abnext
  \begin{tcolorbox}[enhanced, breakable=false, arc=2mm, colback=Paper, colframe=Line,
    left=2mm, right=2mm, top=1mm, bottom=1mm, boxrule=1pt, height=\abh{#1}]
  \end{tcolorbox}\ablog%
}
% \linedbox{height}: an answer box ruled in dotted thirds. Lines are \thirdrow/3 apart and run dotted, dotted,
% solid (the baseline), repeating, and end on one dotted line. The number of
% rows is worked out from the box's final height, so the ruling still fills the
% box after fill.py has grown it; the rows are centred in the box.
\newlength\thirdrow \setlength\thirdrow{9mm}
\tikzset{thirdsolid/.style={draw=DarkGrey!70, line width=0.5pt},
         thirddot/.style={draw=DarkGrey!70, line width=0.6pt, dash pattern=on 0.9pt off 1.6pt}}
\newlength\lbH \newlength\lbO \newcount\lbN
\newcommand{\linedbox}[1]{%
  \abnext
  % interior height = box height less the 1pt frame top and bottom
  \setlength\lbH{\dimexpr\abh{#1}-2pt\relax}%
  % whole rows that fit with at least 3 mm spare top and bottom (at least one)
  \lbN=\numexpr(\dimexpr\lbH-6mm-\thirdrow/2\relax)/\thirdrow\relax
  \ifnum\lbN<1 \lbN=1 \fi
  \setlength\lbO{\dimexpr(\lbH-\lbN\thirdrow)/2\relax}%
  \begin{tcolorbox}[enhanced, breakable=false, arc=2mm, colback=Paper, colframe=Line,
    left=2mm, right=2mm, top=1mm, bottom=1mm, boxrule=1pt, height=\abh{#1},
    underlay={\begin{tcbclipinterior}
      \foreach \i in {1,...,\the\lbN}{
        \draw[thirddot] ([xshift=3mm,yshift=-\lbO-(\i-1)*\thirdrow]interior.north west) -- ([xshift=-3mm,yshift=-\lbO-(\i-1)*\thirdrow]interior.north east);
        \draw[thirddot] ([xshift=3mm,yshift=-\lbO-(\i-1)*\thirdrow-\thirdrow/3]interior.north west) -- ([xshift=-3mm,yshift=-\lbO-(\i-1)*\thirdrow-\thirdrow/3]interior.north east);
        \draw[thirdsolid] ([xshift=3mm,yshift=-\lbO-(\i-1)*\thirdrow-2*\thirdrow/3]interior.north west) -- ([xshift=-3mm,yshift=-\lbO-(\i-1)*\thirdrow-2*\thirdrow/3]interior.north east);}
      \draw[thirddot] ([xshift=3mm,yshift=-\lbO-\lbN\thirdrow]interior.north west) -- ([xshift=-3mm,yshift=-\lbO-\lbN\thirdrow]interior.north east);
    \end{tcbclipinterior}}]
  \end{tcolorbox}\ablog%
}
% Every written-response box is ruled the same way (Part A comprehension, We do,
% Part C tasks and Part D extensions, including side-by-side pairs). Space for
% mathematical working uses \workbox{height} instead and stays plain, as do
% \drawbox and the write-in tables.
\let\workbox\answerbox   % plain box: mathematical working (calculations) stays unruled
\let\answerbox\linedbox
% \drawbox{height}{contents}: white box for a diagram or drawing
\newcommand{\drawbox}[2]{%
  \abnext
  \begin{tcolorbox}[colback=white, colframe=MidGrey, boxrule=0.5pt, height=\abh{#1}]#2\end{tcolorbox}\ablog}
% Notes box at the end of Part B: fills whatever is left of the page (if worth it)
\newcommand{\notesfill}{\par
  \ifdim\dimexpr\pagegoal-\pagetotal\relax>3.2cm
    \vspace*{\fill}%
    \begin{tcolorbox}[enhanced, breakable=false, arc=2mm, colback=Paper, colframe=Line,
      left=2mm, right=2mm, top=1mm, bottom=1mm, boxrule=1pt,
      height=\dimexpr\pagegoal-\pagetotal-1.0cm\relax]
      {\small\headingfont\color{DarkGrey}Notes}
    \end{tcolorbox}%
  \fi}



% Centred table: wrap every content table in \begin{ctab} ... \end{ctab}
% (figures already centre themselves inside their \diag... macros)
% ctab also registers write-in tables (S column) for page filling: the extra
% height is per row, logged with the row count.
\newenvironment{ctab}{\par\smallskip\begingroup\centering
  \abnext\let\tbcur\abcur\global\wtrows0 }{\par
  \ifnum\wtrows>0 \edef\abtmp{\noexpand\write\abpos{\theabox\space\abgrp\space\noexpand\thepage\space\the\wtrows}}\abtmp\fi
  \endgroup\smallskip}

% ── Booklet formatting rules (Oct 2026) ─────────────────────────────────────
% 1. No marks are shown anywhere in the booklet.
\renewcommand{\markscount}[1]{}
% 2. A task's heading sits on its own line; its text starts on the next
%    line, left aligned. Parts (a), (b)… each start a new line: use \pt{a}.
\RenewDocumentCommand{\task}{o}{\par\noindent\refstepcounter{tasknum}\label{task:\thetasknum\IfValueT{#1}{#1}}\textbf{Task \thetasknum\IfValueT{#1}{#1}}\par\nobreak\noindent\ignorespaces}
\renewcommand{\taskpart}[1]{\par\noindent\label{task:\thetasknum#1}\textbf{Task \thetasknum#1}\par\nobreak\noindent\ignorespaces}
\renewcommand{\exttask}{\par\noindent\refstepcounter{tasknum}\label{task:\thetasknum}\textbf{Extension \thetasknum}\par\nobreak\noindent\ignorespaces}
\newcommand{\pt}[1]{\par\noindent\textbf{(#1)}\ \ignorespaces}
% Workstation pages (Experiment 1): numbered method steps with a hanging
% indent, and a one-line safety reminder under the materials list.
\newcommand{\wsstep}[1]{\par\noindent\hangindent=1.6em\hangafter=1\makebox[1.6em][l]{\textbf{#1}}\ignorespaces}
\newcommand{\wssafety}[1]{\par\smallskip{\small\textcolor{ScaffoldRed}{\headingfont\bfseries Safety:}\ #1}\par}
% 3. Tables have vertical rules between columns and down both sides
%    (write | between and around the columns, e.g. {|L{3cm}|C{2cm}|}).
%    Keep the normal padding at the outer edges (no @{} there): colortbl paints
%    header colour one \tabcolsep past each edge, so trimming the padding makes
%    the header overhang the grid.
\arrayrulecolor{MidGrey}
\setlength{\arrayrulewidth}{0.6pt}
% 4. Write-in tables: W/Z are left/centred columns that centre text
%    vertically; put the S column first (see below).
\newcolumntype{W}[1]{>{\RaggedRight\arraybackslash}m{#1}}
\newcolumntype{Z}[1]{>{\Centering\arraybackslash}m{#1}}
\def\tbcur{0pt}
\newcount\wtrows
\newcommand{\wh}{\global\advance\wtrows1 \rule[\dimexpr-0.75cm-\tbcur/2\relax]{0pt}{\dimexpr1.5cm+\tbcur\relax}}
% S = zero-width first column that makes every row of a write-in table
% at least 1.5 cm tall with its text vertically centred. Start each data
% row with "&"; start the header row with "\hd &" so it stays slim.
\newif\ifnowh
\newcolumntype{S}{@{}>{}c<{\ifnowh\else\wh\fi\global\nowhfalse}@{\hspace{\tabcolsep}}}
\newcommand{\hd}{\global\nowhtrue}
% 5. The first row of a multi-row write-in table is a worked scaffold in red.
\definecolor{ScaffoldRed}{HTML}{C0392B}
\newcommand{\scaf}[1]{\textcolor{ScaffoldRed}{\itshape #1}}
% 6. Tables are a light grid (rules between rows and columns) on white —
%    row shading is off, because colortbl paints shading over column rules.
\renewcommand{\talt}{}
\renewcommand{\tmid}{}
% ══════════════════════════════════════════════════════════════════════════════
%  EXPERIMENT BOOKLET LAYER (Nov 2026)
%  From 2027 every Year 7 unit collects its experiments into one A4 experiment
%  booklet that puts the experimental techniques and skills first. Every skill
%  comes from the faculty Science Skills continuum (Years 7--8 column) and is
%  named the same way everywhere: a code (QP1, PC4, PA2 …) on a chip coloured by
%  its sub-strand. Copy this block into the next unit's experiment booklet.
%
%  Each experiment runs in the same order:
%    \experimentdiv{title}{summary}{lesson}\label{sec:N}
%    key vocabulary · learning intentions · \begin{skillsfocus}
%    \exppart*{Technique spotlight}{Blue}   (technique, worked example, We do)
%    \exppart{Plan and safety}{Teal}        (aim, prediction, variables, risks)
%    \exppart{Method}{Teal}
%    \exppart{Results}{Amber}
%    \exppart{Analysis and evaluation}{Coral}
%    \exppart{Skill check}{NavyBlue}
% ══════════════════════════════════════════════════════════════════════════════
\usetikzlibrary{arrows.meta,calc,positioning}
\usepackage{pifont}

\newenvironment{keepext}{\par\begin{tcolorbox}[blankest, breakable=false]}{\end{tcolorbox}\par}
% graph-paper box: \graphgrid{height}
\newcommand{\graphgrid}[1]{%
  \begin{tcolorbox}[enhanced, breakable=false, colback=white, colframe=MidGrey, boxrule=0.5pt, height=#1,
    underlay={\begin{tcbclipinterior}\draw[step=5mm, Line, very thin] (interior.south west) grid (interior.north east);
      \draw[step=1cm, MidGrey!70, thin] (interior.south west) grid (interior.north east);\end{tcbclipinterior}}]
  \end{tcolorbox}}

% ── Skills (faculty Science Skills continuum, Years 7–8) ─────────────────────
% \defskill{code}{name}{colour}{VC code}{what it means}
\makeatletter
\newcommand{\defskill}[5]{%
  \expandafter\gdef\csname sk@n@#1\endcsname{#2}%
  \expandafter\gdef\csname sk@c@#1\endcsname{#3}%
  \expandafter\gdef\csname sk@v@#1\endcsname{#4}%
  \expandafter\gdef\csname sk@d@#1\endcsname{#5}}
\newcommand{\skname}[1]{\csname sk@n@#1\endcsname}
\newcommand{\skcode}[1]{\csname sk@v@#1\endcsname}
\newcommand{\skdesc}[1]{\csname sk@d@#1\endcsname}
% \sk{PC4}: the chip — code and name on the sub-strand colour
\newcommand{\sk}[1]{%
  \edef\sk@col{\csname sk@c@#1\endcsname}%
  \tikz[baseline=(c.base)]\node[rounded corners=2.5pt, fill=\sk@col!16, draw=\sk@col, line width=0.6pt,
    text=\sk@col!55!black, inner xsep=4pt, inner ysep=2pt, font=\headingfont\bfseries\scriptsize]
    (c) {#1\,\textperiodcentered\,\skname{#1}};}
% \skr{PC4}: the chip flush right at the end of a paragraph; drops to its own line if it does not fit
\newcommand{\skr}[1]{\unskip\nobreak\hfill\penalty50\hskip1em\null\nobreak\hfill\sk{#1}}
% \skc{PC4}: the code alone, for tight table cells
\newcommand{\skc}[1]{%
  \edef\sk@col{\csname sk@c@#1\endcsname}%
  \tikz[baseline=(c.base)]\node[rounded corners=2.5pt, fill=\sk@col!16, draw=\sk@col, line width=0.6pt,
    text=\sk@col!55!black, inner xsep=3pt, inner ysep=2pt, font=\headingfont\bfseries\scriptsize]
    (c) {#1};}
\makeatother

% Questioning and predicting — VC2S8I01
\defskill{QP1}{Testable hypothesis}{Teal}{VC2S8I01}{A prediction that an experiment can test, written as \emph{If \dots\ then \dots\ because \dots}}
\defskill{QP2}{Formulate an aim}{Teal}{VC2S8I01}{One sentence saying what the experiment will find out: \emph{To find out how \dots\ affects \dots}}
% Planning and conducting — VC2S8I02, VC2S8I03
\defskill{PC1}{Plan and follow a method}{Blue}{VC2S8I02}{Follow numbered steps in order, so that someone else could repeat the experiment and get the same results (repeatable)}
\defskill{PC2}{Identify variables}{Blue}{VC2S8I02}{Name the independent variable (what you change), the dependent variable (what you measure) and the controlled variables (what you keep the same)}
\defskill{PC3}{Risks and safety}{Blue}{VC2S8I02}{Name each hazard, the risk it causes and the control that prevents it; follow the lab rules}
\defskill{PC4}{Precision}{Blue}{VC2S8I02}{Choose equipment with a fine enough scale, read it to its smallest division, and repeat measurements to see how closely they agree}
\defskill{PC5}{Ethics}{Blue}{VC2S8I02}{Treat people, living things, the environment and other people's data with respect}
\defskill{PC6}{Collect data and measure}{Blue}{VC2S8I02}{Take each measurement at the right time, with its unit, and record it straight away}
\defskill{PC7}{Use and select equipment}{Blue}{VC2S8I03}{Name each piece of equipment, use it correctly and choose the right piece for the job}
\defskill{PC8}{Use digital tools}{Blue}{VC2S8I03}{Use stopwatches, data loggers, sensors and spreadsheets to collect and record data}
% Processing, modelling and analysing — VC2S8I04, VC2S8I05
\defskill{PA1}{Tables}{GoldDark}{VC2S8I04}{Organise data in a ruled table, with the independent variable first and units in the column headings}
\defskill{PA2}{Graphing}{GoldDark}{VC2S8I04}{Draw the right type of graph with a title, labelled axes with units, even scales and accurate points}
\defskill{PA3}{Mathematical relationships}{GoldDark}{VC2S8I04}{Process data with calculations: differences, averages, rates and formulas}
\defskill{PA4}{Creating models}{GoldDark}{VC2S8I04}{Build a physical, drawn or mathematical model that represents something real}
\defskill{PA5}{Patterns and trends}{GoldDark}{VC2S8I05}{Describe how the dependent variable changes as the independent variable changes, quoting data}
\defskill{PA6}{Correlation and causality}{GoldDark}{VC2S8I05}{Say whether two variables change together, and whether there is evidence that one causes the other}
\defskill{PA7}{Identify anomalies}{GoldDark}{VC2S8I05}{Spot a result that does not fit the pattern and suggest why it happened}
\defskill{PA8}{Interpreting models}{GoldDark}{VC2S8I05}{Use a model to explain an observation, and say what the model shows well and what it leaves out}
% Evaluating — VC2S8I06, VC2S8I07
\defskill{EV1}{Results against hypothesis}{OrangeDark}{VC2S8I06}{Say whether the results support the hypothesis, quoting data as evidence}
\defskill{EV2}{Compare with known concepts}{OrangeDark}{VC2S8I06}{Compare your results with accepted values and scientific ideas}
\defskill{EV3}{Error sources}{OrangeDark}{VC2S8I06}{Name what could have made results less accurate or less precise, and suggest how to improve the method}
\defskill{EV4}{Quality of secondary sources}{OrangeDark}{VC2S8I07}{Judge whether information from a website, book or video can be trusted}
\defskill{EV5}{Questions for further investigation}{OrangeDark}{VC2S8I07}{Pose a new question that follows from your results}
% Communicating — VC2S8I08
\defskill{CO1}{Scientific writing}{PurpleDark}{VC2S8I08}{Write in the third person, legibly, with correct spelling and grammar, short and to the point, using scientific terms}
\defskill{CO2}{Presentation formats}{PurpleDark}{VC2S8I08}{Choose the right format to show findings: a table, graph, scientific diagram, poster, report or talk}

% Skills focus panel at the top of each experiment:
%   \begin{skillsfocus} \focus{PC4}{what you do with it here} … \end{skillsfocus}
\newtcolorbox{skillsfocus}{
  enhanced, breakable=false, arc=2.5mm, boxrule=0.6pt,
  colback=white, colframe=NavyBlue!55,
  fonttitle=\headingfont\bfseries, coltitle=white, colbacktitle=NavyBlue,
  title={Skills focus --- the science skills this experiment builds},
  left=4mm, right=3.5mm, top=2mm, bottom=2mm}
\newcommand{\focus}[2]{\par\noindent\makebox[1.25cm][l]{\skc{#1}}\parbox[t]{\dimexpr\linewidth-1.25cm\relax}{\small\textbf{\skname{#1}}\,---\,#2}\par\vspace{2pt}}

% Learning intentions panel: \begin{lipanel} \inquiryli{…} \understandli{…} \sclist{…} \end{lipanel}
\newtcolorbox{lipanel}{
  enhanced, breakable=false, arc=2.5mm, boxrule=0.5pt,
  colback=LightBlue, colframe=Blue!40,
  fonttitle=\headingfont\bfseries, coltitle=MidBlue, title={Learning intentions},
  left=4mm, right=3.5mm, top=2mm, bottom=2mm}
\newcommand{\inquiryli}[1]{\par\noindent{\headingfont\bfseries\color{NavyBlue}Science Inquiry:}\ #1\par\vspace{2pt}}
\newcommand{\understandli}[1]{\par\noindent{\headingfont\bfseries\color{NavyBlue}Science Understanding:}\ #1\par\vspace{2pt}}
\newcommand{\sctitle}{\par\noindent{\headingfont\bfseries\color{NavyBlue}Success criteria}\par}

% Technique spotlight: \begin{technique}{title}{PC4} … \end{technique}
\newtcolorbox{technique}[2]{
  enhanced, breakable, arc=2.5mm, boxrule=0.6pt,
  colback=white, colframe=Blue!55,
  fonttitle=\headingfont\bfseries, coltitle=MidBlue, colbacktitle=LightBlue,
  title={Technique: #1\hfill\sk{#2}},
  left=4mm, right=3.5mm, top=2.5mm, bottom=2.5mm}

% Experiment divider: \experimentdiv{title}{summary}{lesson number}
% Steps the section counter, so tasks, worked examples and figures are N.M.
\newcommand{\experimentdiv}[3]{%
  \clearpage
  \refstepcounter{section}%
  \pdfbookmark[1]{Experiment \thesection{} — #1}{exp.\thesection}%
  \begin{tcolorbox}[
    enhanced, breakable=false, arc=3.5mm,
    colback=NavyBlue, colframe=NavyBlue, boxrule=0pt,
    sidebyside, sidebyside align=center,
    lefthand width=2.6cm,
    left=2mm, right=4mm, top=4mm, bottom=4mm]
    \begin{center}
      {\headingfont\bfseries\color{white!60!NavyBlue}\scriptsize EXPERIMENT}\\[-2pt]
      {\headingfont\bfseries\color{white!35!NavyBlue}\fontsize{46}{46}\selectfont \thesection}
    \end{center}
    \tcblower
    {\headingfont\bfseries\color{white}\LARGE #1}\\[3pt]
    {\itshape\color{white!75!NavyBlue}\small #2}\\[3pt]
    {\headingfont\color{Sky}\footnotesize Lesson #3}
  \end{tcolorbox}%
}

% Part bar inside an experiment. \exppart{title}{colour} starts a new page
% (new parts always start on a new page); \exppart*{…}{…} follows straight on.
\NewDocumentCommand{\exppart}{s m m}{%
  \par\IfBooleanTF{#1}{\needspace{10\baselineskip}\vspace{8pt}}{\clearpage}%
  \begin{tcolorbox}[enhanced, breakable=false, arc=2mm, boxrule=0pt,
    colback=#3!14, borderline west={4pt}{0pt}{#3},
    left=4mm, right=3mm, top=1.6mm, bottom=1.6mm]
  {\headingfont\bfseries\color{#3!60!black}\large #2}
  \end{tcolorbox}\vspace{2pt}}

% Skill check at the end of each experiment: students tick one column per row.
%   \begin{skillcheck} \scrow{PC4}{I can …} … \end{skillcheck}
\newcommand{\tickbox}{\tikz[baseline=-0.6ex]\draw[DarkGrey, line width=0.7pt, rounded corners=1pt] (0,-0.17) rectangle (0.34,0.17);}
\newenvironment{skillcheck}{%
  \par\noindent Read each statement. Tick the column that best describes you \emph{today}.\par
  \begin{ctab}\renewcommand{\arraystretch}{1.25}%
  \begin{tabular}{|W{3.9cm}|W{7.0cm}|Z{1.45cm}|Z{1.45cm}|Z{1.45cm}|}
  \hline\thead{Skill} & \thead{Success criterion} & \thead{Not yet} & \thead{Nearly} & \thead{Got it} \\ \hline}
  {\end{tabular}\end{ctab}}
\newcommand{\scrow}[2]{\rule{0pt}{1.1cm}\skc{#1}\enspace{\small\skname{#1}} & #2 & \tickbox & \tickbox & \tickbox \\ \hline}

% Write-in label under a drawing: \nameline{label}
\newcommand{\nameline}[1]{\par\noindent{\small #1}\ \leaders\hbox{\rule[-2pt]{1pt}{0.4pt}}\hfill\null\par}

% ── Equipment drawn as scientific diagrams (2D, side-on, outline only) ──────
% Each \eq... macro draws at the origin inside a tikzpicture. \eqpic[scale]{\eqBeaker}
% puts one on its own. The same shapes are used everywhere, so the booklet
% models exactly the drawings students are asked to make.
\tikzset{
  eq/.style={line width=0.9pt, draw=Ink, line cap=round, line join=round},
  eqlab/.style={draw=Ink, line width=0.4pt},
  eqtext/.style={font=\footnotesize, text=Ink, inner sep=1.5pt},
}
\newcommand{\eqpic}[2][1]{\begin{tikzpicture}[scale=#1, baseline=(current bounding box.center)]#2\end{tikzpicture}}
\newcommand{\eqBeaker}{\draw[eq] (-0.22,2.62) -- (0,2.45) -- (0,0) -- (1.8,0) -- (1.8,2.5);}
\newcommand{\eqFlask}{\draw[eq] (0.62,2.9) -- (0.65,2.8) -- (0.65,1.9) -- (0,0) -- (1.9,0) -- (1.25,1.9) -- (1.25,2.8) -- (1.28,2.9);}
\newcommand{\eqTestTube}{\draw[eq] (-0.06,3.0) -- (0,2.9) -- (0,0.3) arc[start angle=180,end angle=360,radius=0.3] -- (0.6,2.9) -- (0.66,3.0);}
\newcommand{\eqCylinder}{%
  \draw[eq] (-0.15,3.6) -- (0,3.5) -- (0,0.15) (0.6,0.15) -- (0.6,3.5);
  \draw[eq] (-0.35,0) rectangle (0.95,0.15);
  \foreach \y in {0.6,1.1,...,3.2} \draw[Ink, line width=0.5pt] (0.6,\y) -- (0.4,\y);}
\newcommand{\eqFunnel}{\draw[eq] (0,2.2) -- (0.75,1.0) -- (0.75,0) (1.85,2.2) -- (1.1,1.0) -- (1.1,0);}
\newcommand{\eqBasin}{\draw[eq] (0,0.9) .. controls (0.15,0.02) and (2.25,0.02) .. (2.4,0.9);}
\newcommand{\eqTripod}{%
  \draw[eq] (-0.15,2.0) rectangle (2.55,2.12);
  \draw[eq] (0.25,2.0) -- (0,0) (2.15,2.0) -- (2.4,0) (1.2,2.0) -- (1.2,0.15);}
\newcommand{\eqBunsen}{%
  \draw[eq] (0,0) rectangle (1.6,0.22);
  \draw[eq] (0.58,0.85) -- (0.58,2.3) (1.02,0.85) -- (1.02,2.3);
  \draw[eq] (0.5,0.22) rectangle (1.1,0.85);
  \fill[Ink] (0.72,0.42) rectangle (0.88,0.62);
  \draw[eq] (-0.65,0.06) rectangle (0,0.16);}
\newcommand{\eqThermometer}{%
  \draw[eq] (0.25,0.25) circle (0.2);
  \draw[eq] (0.15,0.43) -- (0.15,3.9) arc[start angle=180,end angle=0,radius=0.1] -- (0.35,0.43);
  \draw[Ink, line width=1.3pt] (0.25,0.45) -- (0.25,1.9);}
\newcommand{\eqStand}{%
  \draw[eq] (0,0) rectangle (2.3,0.2);
  \draw[eq] (0.35,0.2) -- (0.35,4.0) (0.5,0.2) -- (0.5,4.0);
  \draw[eq] (0.25,2.75) rectangle (0.75,3.2);
  \draw[eq] (0.75,2.92) -- (1.6,2.92) (0.75,3.03) -- (1.6,3.03);
  \draw[eq] (1.6,2.7) -- (1.6,3.25) (1.6,2.7) -- (1.95,2.7) (1.6,3.25) -- (1.95,3.25);}
\newcommand{\eqWatchGlass}{\draw[eq] (0,0.35) .. controls (0.5,-0.05) and (1.7,-0.05) .. (2.2,0.35);}
\newcommand{\eqTongs}{%
  \draw[eq] (0,0.1) -- (2.0,0.45) .. controls (2.3,0.5) and (2.45,0.35) .. (2.6,0.15)
            (0,0.0) -- (2.0,-0.35) .. controls (2.3,-0.4) and (2.45,-0.25) .. (2.6,-0.05);}

% A label line from a point on a drawing to text at the side: \eqlab{(x,y)}{(tx,ty)}{text}
\newcommand{\eqlab}[3]{\draw[eqlab] #1 -- #2 node[eqtext, anchor=west] {#3};}
\newcommand{\eqlabl}[3]{\draw[eqlab] #1 -- #2 node[eqtext, anchor=east] {#3};}

% ── Reading scales ──────────────────────────────────────────────────────────
% \readcyl{lowest}{highest}{labelled step}{small divisions per label}{reading}
% A close-up of a measuring cylinder's scale, the water drawn to the reading:
% the bottom of the meniscus sits exactly on the reading.
\pgfkeys{/pgf/number format/.cd, fixed, precision=1}
\newcommand{\readcyl}[5]{%
\begin{tikzpicture}[baseline=(current bounding box.south)]
  \pgfmathsetmacro\rcH{5.0}
  \pgfmathsetmacro\rcY{(#5-#1)*\rcH/(#2-#1)}
  \pgfmathtruncatemacro\rcN{round((#2-#1)/#3*#4)}
  \fill[black!10] (0,-0.3) -- (0,\rcY+0.16) .. controls (0.45,\rcY-0.05) and (1.15,\rcY-0.05) .. (1.6,\rcY+0.16) -- (1.6,-0.3) -- cycle;
  \draw[Ink, line width=0.9pt] (0,\rcY+0.16) .. controls (0.45,\rcY-0.05) and (1.15,\rcY-0.05) .. (1.6,\rcY+0.16);
  \foreach \k in {0,...,\rcN}{
    \pgfmathsetmacro\yy{\k*\rcH/\rcN}
    \pgfmathtruncatemacro\ismaj{mod(\k,#4)==0 ? 1 : 0}
    \ifnum\ismaj=1
      \draw[Ink, line width=0.75pt] (0,\yy) -- (0.6,\yy);
      \pgfmathsetmacro\lv{#1+\k/#4*#3}
      \node[anchor=west, font=\scriptsize\bfseries, inner sep=1pt, fill=white, fill opacity=0.8, text opacity=1] at (0.66,\yy) {\pgfmathprintnumber{\lv}};
    \else
      \draw[Ink, line width=0.4pt] (0,\yy) -- (0.32,\yy);
    \fi}
  \draw[eq] (0,-0.3) -- (0,\rcH+0.35) (1.6,-0.3) -- (1.6,\rcH+0.35);
  \node[font=\scriptsize, anchor=north] at (0.8,-0.35) {scale in mL};
\end{tikzpicture}}
% \readtherm{lowest}{highest}{labelled step}{small divisions per label}{reading}
\newcommand{\readtherm}[5]{%
\begin{tikzpicture}[baseline=(current bounding box.south)]
  \pgfmathsetmacro\rtH{5.0}
  \pgfmathsetmacro\rtY{(#5-#1)*\rtH/(#2-#1)}
  \pgfmathtruncatemacro\rtN{round((#2-#1)/#3*#4)}
  \draw[eq, fill=white] (0,-0.45) circle (0.3);
  \fill[black!75] (0,-0.45) circle (0.22);
  \fill[black!75] (-0.07,-0.3) rectangle (0.07,\rtY);
  \draw[eq] (-0.16,-0.2) -- (-0.16,\rtH+0.3) arc[start angle=180,end angle=0,radius=0.16] -- (0.16,-0.2);
  \foreach \k in {0,...,\rtN}{
    \pgfmathsetmacro\yy{\k*\rtH/\rtN}
    \pgfmathtruncatemacro\ismaj{mod(\k,#4)==0 ? 1 : 0}
    \ifnum\ismaj=1
      \draw[Ink, line width=0.75pt] (0.25,\yy) -- (0.75,\yy);
      \pgfmathsetmacro\lv{#1+\k/#4*#3}
      \node[anchor=west, font=\scriptsize\bfseries, inner sep=1pt] at (0.8,\yy) {\pgfmathprintnumber{\lv}};
    \else
      \draw[Ink, line width=0.4pt] (0.25,\yy) -- (0.52,\yy);
    \fi}
  \node[font=\scriptsize, anchor=north] at (0.3,-0.8) {scale in \textdegree C};
\end{tikzpicture}}
% \balancefig{reading}: an electronic balance showing a reading
\newcommand{\balancefig}[1]{%
\begin{tikzpicture}[baseline=(current bounding box.south)]
  \draw[eq, rounded corners=3pt] (0,0) rectangle (4.2,1.35);
  \draw[eq] (1.9,1.35) -- (1.9,1.55) (2.3,1.35) -- (2.3,1.55);
  \draw[eq, rounded corners=1pt] (0.5,1.55) rectangle (3.7,1.72);
  \draw[eq, fill=LightGrey] (0.35,0.35) rectangle (2.75,1.02);
  \node[font=\headingfont\bfseries\large, anchor=east] at (2.65,0.68) {#1};
  \draw[eq] (3.3,0.68) circle (0.2);
  \node[font=\tiny\bfseries] at (3.3,0.3) {TARE};
\end{tikzpicture}}
% Bunsen burner with a flame: \bunsenflame{safety|heating}
\newcommand{\bunsenflame}[1]{%
\begin{tikzpicture}[baseline=(current bounding box.south)]
  \def\fl{#1}\def\safe{safety}%
  \ifx\fl\safe
    \fill[black!14] (0.8,2.3) .. controls (0.05,3.0) and (0.45,4.3) .. (0.85,5.0) .. controls (1.05,4.4) and (1.55,3.1) .. (0.8,2.3);
    \draw[Ink, line width=0.6pt] (0.8,2.3) .. controls (0.05,3.0) and (0.45,4.3) .. (0.85,5.0) .. controls (1.05,4.4) and (1.55,3.1) .. (0.8,2.3);
  \else
    \fill[white] (0.8,2.3) .. controls (0.45,2.8) and (0.6,3.5) .. (0.8,3.9) .. controls (1.0,3.5) and (1.15,2.8) .. (0.8,2.3);
    \draw[Ink, line width=0.6pt] (0.8,2.3) .. controls (0.45,2.8) and (0.6,3.5) .. (0.8,3.9) .. controls (1.0,3.5) and (1.15,2.8) .. (0.8,2.3);
    \fill[black!30] (0.8,2.3) .. controls (0.65,2.55) and (0.7,2.85) .. (0.8,3.0) .. controls (0.9,2.85) and (0.95,2.55) .. (0.8,2.3);
  \fi
  \draw[eq] (0,0) rectangle (1.6,0.22);
  \draw[eq] (0.58,0.85) -- (0.58,2.3) (1.02,0.85) -- (1.02,2.3);
  \draw[eq] (0.5,0.22) rectangle (1.1,0.85);
  \ifx\fl\safe\draw[Ink, line width=0.5pt, densely dotted] (0.72,0.42) rectangle (0.88,0.62);\else\fill[Ink] (0.72,0.42) rectangle (0.88,0.62);\fi
  \draw[eq] (-0.65,0.06) rectangle (0,0.16);
\end{tikzpicture}}
% ══════════════════════════════════════════════════════════════════════════════
%  SEPARATION — COVER
%  The unit's cover art (covers/separation.jpg), full bleed. The A5 lesson
%  booklet sets \stcoverfile one folder deeper.
% ══════════════════════════════════════════════════════════════════════════════
\providecommand{\stcoverfile}{../../covers/separation.jpg}
\newcommand{\stcover}{%
\begin{tikzpicture}[remember picture, overlay]
  \node[anchor=north] at (current page.north) {\includegraphics[width=\paperwidth]{\stcoverfile}};
\end{tikzpicture}}
% ══════════════════════════════════════════════════════════════════════════════
%  2027 WORKBOOK STYLE (Year 7 workbook redesign, Nov 2026) — A4 version
%  The look of the A5 trial (space/booklet/a5-test) on an A4 page:
%    · colour cover from the unit's cover art, black-and-white inside
%    · headings in Outfit (the site's heading face), body in Lato
%    · each experiment opens with a band: number, title, ink symbol, heavy rule
%    · boxes told apart by frame and icon, never by colour:
%        I do ●  solid · We do ●● dashed · You do ●●● heavy
%        Key idea / technique ◆ thin · Above level ▲ double · Country ◎ dotted
%    · skill tags are ink tabs carrying the code (QP1, PC4 …) plus the name
%  Everything above this block still names the old colours; they are mapped to
%  ink here, so the body never needs to know which look it is printed in.
% ══════════════════════════════════════════════════════════════════════════════

\usetikzlibrary{patterns}
% ── Fonts ───────────────────────────────────────────────────────────────────
% Bold is Outfit SemiBold, the weight the site uses; Outfit has no italic,
% so a slanted one stands in.
\renewfontfamily\headingfont{Outfit}[
  Path=./fonts/, Extension=.ttf, Ligatures=TeX,
  UprightFont=Outfit-Regular, BoldFont=Outfit-SemiBold,
  ItalicFont=Outfit-Regular, ItalicFeatures={FakeSlant=0.18},
  BoldItalicFont=Outfit-SemiBold, BoldItalicFeatures={FakeSlant=0.18} ]

% ── Black and white inside ──────────────────────────────────────────────────
\definecolor{Ink}{HTML}{111111}
\colorlet{NavyBlue}{Ink}\colorlet{MidBlue}{Ink}\colorlet{TealGreen}{Ink}
\colorlet{GoldDark}{Ink}\colorlet{PurpleDark}{Ink}\colorlet{OrangeDark}{Ink}
\colorlet{TableHead}{Ink}
\colorlet{Blue}{Ink}\colorlet{Teal}{Ink}\colorlet{Amber}{Ink}\colorlet{Coral}{Ink}
\colorlet{Violet}{Ink}\colorlet{Sky}{Ink}
\colorlet{LightBlue}{white}\colorlet{LightGreen}{white}\colorlet{GoldLight}{white}
\colorlet{PurpleLight}{white}\colorlet{OrangeLight}{white}\colorlet{RowAlt}{white}
\colorlet{LightGrey}{white}\colorlet{Paper}{white}
\colorlet{DarkGrey}{black!72}\colorlet{MidGrey}{black!45}\colorlet{Line}{black!45}
\arrayrulecolor{black!60}

% ── Page style ──────────────────────────────────────────────────────────────
\fancypagestyle{fancy}{%
  \fancyhf{}%
  \renewcommand{\headrulewidth}{0.4pt}%
  \renewcommand{\headrule}{\hbox to\headwidth{\color{Ink}\leaders\hrule height \headrulewidth\hfill}}%
  \fancyhead[L]{\headingfont\footnotesize \unitname\ \textbullet\ Experiment Booklet}%
  \fancyhead[R]{\headingfont\footnotesize\leftmark}%
  \fancyfoot[C]{\headingfont\footnotesize\thepage}}
\fancypagestyle{plain}{\fancyhf{}\renewcommand{\headrulewidth}{0pt}%
  \fancyfoot[C]{\headingfont\footnotesize\thepage}}
\fancypagestyle{opener}{\fancyhf{}\renewcommand{\headrulewidth}{0pt}%
  \fancyfoot[C]{\headingfont\footnotesize\thepage}}
\pagestyle{fancy}
\markboth{}{}

\titleformat{\subsubsection}{\headingfont\normalsize\bfseries}{}{0em}{}

% ── Box system: told apart by frame and icon, never by colour ───────────────
\newcommand{\icodots}[1]{\tikz[baseline=-0.6ex]{\foreach \k in {1,...,#1}{\fill[white] (\k*1.9mm,0) circle (0.7mm);}}}
\newcommand{\icokey}{\tikz[baseline=-0.6ex]{\fill[white] (0,0.9mm)--(0.9mm,0)--(0,-0.9mm)--(-0.9mm,0)--cycle;}}
\newcommand{\icoup}{\tikz[baseline=-0.6ex]{\fill[white] (-1mm,-0.8mm)--(1mm,-0.8mm)--(0,1mm)--cycle;}}
\newcommand{\icocountry}{\tikz[baseline=-0.6ex]{\draw[white,line width=0.5pt] (0,0) circle (1mm); \fill[white] (0,0) circle (0.45mm);}}
\newcommand{\icoskill}{\tikz[baseline=-0.6ex]{\fill[white] (-0.85mm,-0.85mm) rectangle (0.85mm,0.85mm);}}
\newcommand{\tagicon}[1]{\tikz[baseline=-0.6ex]\node[fill=Ink, rounded corners=0.8mm, inner xsep=1.2mm, inner ysep=1mm]{#1};}
\tcbset{wbbase/.style={enhanced, colback=white, arc=1.2mm,
    left=3.5mm, right=3.5mm, top=4.2mm, bottom=2.6mm, before skip=11pt, after skip=8pt,
    attach boxed title to top left={xshift=3mm, yshift=-2.4mm},
    boxed title style={colback=Ink, colframe=Ink, arc=0.8mm, boxrule=0pt,
                       left=1.6mm, right=1.8mm, top=0.6mm, bottom=0.6mm},
    fonttitle=\headingfont\bfseries\footnotesize, coltitle=white}}

% I do: the worked example (solid frame, one dot)
\renewtcolorbox{workedexbox}[1]{wbbase, breakable, colframe=Ink, boxrule=1pt,
  title={\icodots{1}\ I do \textbullet\ #1}}
\renewenvironment{workedex}[1]
  {\refstepcounter{wenum}\label{we:\thewenum}\begin{workedexbox}{Worked example \thewenum\ --- #1}}
  {\end{workedexbox}}
% We do: straight after its worked example (dashed frame, two dots)
\renewtcolorbox{wedobox}[1]{wbbase, breakable=false, frame hidden,
  borderline={1.1pt}{0pt}{Ink, dash pattern=on 4pt off 2.6pt},
  title={\icodots{2}\ #1}}
\renewenvironment{wedo}[1]
  {\par\begin{wedobox}{We do \thewenum\ --- #1}\label{wd:\thewenum}\ignorespaces}
  {\end{wedobox}}
% You do: a set of practice tasks on your own (heavy frame, three dots)
\newtcolorbox{youdobox}[1][Check your understanding]{wbbase, breakable=false,
  colframe=Ink, boxrule=2.4pt, top=5mm,
  title={\icodots{3}\ You do \textbullet\ #1}}
% Key idea, technique, safety, word banks (thin frame, diamond)
\renewtcolorbox{infobox}[1]{wbbase, breakable, colframe=Ink, boxrule=0.5pt,
  title={\icokey\ #1}}
\renewtcolorbox{scaffold}[1]{wbbase, breakable, colframe=Ink, boxrule=0.5pt,
  title={\icokey\ #1}}
\renewtcolorbox{technique}[2]{wbbase, breakable, colframe=Ink, boxrule=0.5pt,
  title={\icokey\ Technique \textbullet\ #1\enspace\textbar\enspace #2 \skname{#2}}}
% Above level (double frame, triangle)
\renewtcolorbox{yr8box}{wbbase, breakable, frame hidden,
  borderline={0.6pt}{0pt}{Ink}, borderline={0.6pt}{1.6pt}{Ink},
  title={\icoup\ Going further \textbullet\ Above level}}
% Country (dotted frame, ring)
\newtcolorbox{countrybox}[1]{wbbase, breakable, frame hidden,
  borderline={1.4pt}{0pt}{Ink, line cap=round, dash pattern=on 0pt off 2.4pt},
  title={\icocountry\ #1}}
% Learning intentions and skills focus (thin frame; square for skills)
\renewtcolorbox{lipanel}{wbbase, breakable=false, colframe=Ink, boxrule=0.5pt,
  title={\icokey\ Learning intentions}}
\renewtcolorbox{skillsfocus}{wbbase, breakable=false, colframe=Ink, boxrule=1pt,
  title={\icoskill\ Skills focus \textbullet\ the science skills this experiment builds},
  before upper={\begin{multicols}{2}\raggedcolumns}, after upper={\end{multicols}}}

% ── Skill tags: an ink tab with the code, then the name ─────────────────────
\makeatletter
\renewcommand{\skc}[1]{%
  \tikz[baseline=(c.base)]\node[rounded corners=0.8mm, fill=Ink, text=white, inner xsep=1.5mm, inner ysep=0.7mm,
    font=\headingfont\bfseries\scriptsize] (c) {#1};}
\renewcommand{\sk}[1]{\mbox{\skc{#1}\,{\headingfont\bfseries\footnotesize\skname{#1}}}}
\makeatother
\renewcommand{\focus}[2]{\par\noindent\makebox[1.1cm][l]{\skc{#1}}\parbox[t]{\dimexpr\linewidth-1.1cm\relax}{\footnotesize\RaggedRight\textbf{\skname{#1}}\,---\,#2}\par\vspace{3pt}}

% ── Experiment opener band ──────────────────────────────────────────────────
% \experimentdiv{title}{summary}{lesson}: number, title and the experiment's ink
% symbol (\expsymbol{N}), then a heavy rule. Steps the section counter, so
% tasks, worked examples and figures are N.M.
\newcommand{\expword}{Experiment}\newcommand{\expcaps}{EXPERIMENT}\newcommand{\symid}{\thesection}
\renewcommand{\experimentdiv}[3]{%
  \clearpage\thispagestyle{opener}%
  \refstepcounter{section}\markboth{\expword\ \thesection\ \textbullet\ #1}{}%
  \pdfbookmark[1]{\expword\ \thesection{} — #1}{exp.\thesection}%
  \noindent\begin{minipage}[b]{\dimexpr\linewidth-46mm\relax}
    {\headingfont\bfseries\footnotesize\addfontfeatures{LetterSpace=14}\unitcaps\ \textbullet\ \expcaps}\par\vspace{1mm}
    {\headingfont\bfseries\fontsize{40}{40}\selectfont\thesection}\par\vspace{1mm}
    {\headingfont\bfseries\fontsize{19}{22}\selectfont\RaggedRight #1\par}\vspace{1mm}
    {\small\itshape #2}\par\vspace{1mm}
    {\headingfont\footnotesize #3}
  \end{minipage}\hfill
  \begin{minipage}[b]{44mm}\centering\resizebox{!}{27mm}{\expsymbol{\symid}}\end{minipage}\par
  \vspace{2mm}\noindent\rule{\linewidth}{1.6pt}\par\vspace{1mm}}
\renewcommand{\lessoncodes}[1]{{\small\itshape VC 2.0: #1}\par\vspace{2pt}}

% Part heading inside an experiment: every part starts a new page; the
% starred form follows straight on.
\RenewDocumentCommand{\exppart}{s m m}{%
  \par\IfBooleanTF{#1}{\needspace{10\baselineskip}\vspace{8pt}}{\clearpage}%
  {\headingfont\bfseries\LARGE #2}\par\vspace{1mm}\noindent\rule{\linewidth}{0.6pt}\par\vspace{2mm}}

% Key vocabulary: a heading like the trial's, then the write-in table
\renewcommand{\vocabhead}{%
  \par\vspace{4pt}{\headingfont\bfseries\Large Key vocabulary}\par\vspace{1pt}%
  \noindent{\small Copy the meaning of each word from the slide. The first one has been done for you.}\par}

% Table headings: bold Outfit on white, the grid in ink
\renewcommand{\thead}[1]{\headingfont\bfseries #1}
\renewcommand{\tickbox}{\tikz[baseline=-0.6ex]\draw[Ink, line width=0.7pt, rounded corners=1pt] (0,-0.17) rectangle (0.34,0.17);}

% Glossary rows: \glossrow{word}{label}
\newcommand{\glossrow}[2]{\rule{0pt}{5mm}\textbf{#1} & \pageref{#2} \\ \hline}

% ── Ink symbols for the experiment openers ───────────────────────────────────
% One diagrammatic idea per experiment, drawn in ink from the same 2D
% equipment shapes the booklet teaches.
\tikzset{sym/.style={line width=1.5pt, draw=Ink, line cap=round, line join=round}}
\newcommand{\symtube}[3]{% x, fill level, fill
  \fill[#3] (#1,0.3) arc[start angle=180,end angle=360,radius=0.3] -- (#1+0.6,#2) -- (#1,#2) -- cycle;
  \draw[sym] (#1,2.6) -- (#1,0.3) arc[start angle=180,end angle=360,radius=0.3] -- (#1+0.6,2.6);}
\newcommand{\symbeaker}{\draw[sym] (-0.15,2.7) -- (0,2.55) -- (0,0) -- (2.6,0) -- (2.6,2.6);}
\newcommand{\expsymbol}[1]{%
  \begin{tikzpicture}[x=1cm,y=1cm]
  \ifcase#1
  \or % 1 Solubility: three test tubes, one clear, one cloudy, one with a sediment
    \symtube{0}{1.6}{white}\symtube{1.2}{1.6}{black!25}\symtube{2.4}{1.6}{white}
    \fill[black!80] (2.4,0.3) arc[start angle=180,end angle=360,radius=0.3] -- (3.0,0.55) -- (2.4,0.55) -- cycle;
    \foreach \t in {0,1.2,2.4}{\draw[Ink, line width=1pt] (\t,1.6) -- (\t+0.6,1.6);}
  \or % 2 Surface area: a sugar cube beside the same sugar as fine grains
    \path (0,-0.7) -- (0,2.6);
    \draw[sym, fill=white] (0,0) rectangle (1.4,1.4);
    \draw[sym] (0,1.4) -- (0.4,1.8) -- (1.8,1.8) -- (1.4,1.4) (1.8,1.8) -- (1.8,0.4) -- (1.4,0);
    \foreach \x/\y in {2.5/0.1,2.8/0.1,3.1/0.1,3.4/0.1,3.7/0.1,2.65/0.35,2.95/0.35,3.25/0.35,3.55/0.35,2.8/0.6,3.1/0.6,3.4/0.6,2.95/0.85,3.25/0.85,3.1/1.1}{\draw[Ink, line width=1pt, fill=white] (\x,\y) rectangle (\x+0.22,\y+0.22);}
  \or % 3 Fast and slow evaporation: small crystals and large crystals in two dishes
    \draw[sym] (0,0.8) .. controls (0.1,0) and (1.7,0) .. (1.8,0.8);
    \draw[sym] (2.3,0.8) .. controls (2.4,0) and (4.0,0) .. (4.1,0.8);
    \foreach \x in {0.4,0.6,0.8,1.0,1.2,1.4}{\fill[Ink] (\x,0.35) rectangle (\x+0.1,0.45);}
    \foreach \x in {2.75,3.35}{\draw[Ink, line width=1.2pt, fill=white] (\x,0.25) rectangle (\x+0.45,0.7);}
    \draw[Ink, line width=1.2pt] (0.9,1.2) .. controls (0.6,1.6) and (1.2,1.8) .. (0.9,2.3);
    \node[font=\headingfont\bfseries\small] at (0.9,2.6) {fast};
    \node[font=\headingfont\bfseries\small] at (3.2,1.5) {slow};
  \or % 4 Rates of evaporation: a dish on a balance, mass going down
    \draw[sym] (0,0) rectangle (2.6,0.6);
    \draw[sym] (0.3,0.6) -- (0.3,0.8) -- (2.3,0.8) -- (2.3,0.6);
    \draw[sym, fill=white] (0.5,0.2) rectangle (1.7,0.45);
    \draw[sym] (0.7,0.8) .. controls (0.8,1.3) and (1.8,1.3) .. (1.9,0.8);
    \foreach \x in {1.0,1.3,1.6}{\draw[Ink, line width=1.2pt] (\x,1.4) .. controls (\x-0.15,1.65) and (\x+0.15,1.85) .. (\x,2.1);}
    \draw[Ink, line width=2.4pt, -{Stealth[length=8pt]}] (3.2,2.0) -- (3.2,0.4);
  \or % 5 Solubility and temperature: a thermometer and a rising curve
    \draw[sym, -{Stealth[length=7pt]}] (0,0) -- (3.4,0); \draw[sym, -{Stealth[length=7pt]}] (0,0) -- (0,3.0);
    \draw[Ink, line width=2.6pt, line cap=round] (0.2,0.3) .. controls (1.4,0.5) and (2.3,1.1) .. (3.0,2.6);
    \draw[sym, fill=white] (3.9,0.35) circle (0.25);
    \draw[sym] (3.78,0.58) -- (3.78,2.9) arc[start angle=180,end angle=0,radius=0.12] -- (4.02,0.58);
    \draw[Ink, line width=2.4pt] (3.9,0.5) -- (3.9,2.2);
  \or % 6 Copper sulfate crystals: a seed crystal on a thread in a beaker
    \symbeaker
    \draw[Ink, line width=1pt] (0.05,1.9) -- (2.55,1.9);
    \draw[sym] (0.6,2.9) -- (2.0,2.9);
    \draw[Ink, line width=1pt] (1.3,2.9) -- (1.3,1.2);
    \draw[sym, fill=black!60] (1.3,1.25) -- (1.6,0.95) -- (1.45,0.55) -- (1.15,0.55) -- (1.0,0.95) -- cycle;
  \or % 7 Filtration: funnel with filter paper in a conical flask
    \draw[sym] (0,3.2) -- (1.1,1.9) -- (1.1,1.1) (2.6,3.2) -- (1.5,1.9) -- (1.5,1.1);
    \draw[Ink, line width=1.2pt, fill=black!15] (0.3,3.0) -- (1.3,2.0) -- (2.3,3.0) -- cycle;
    \draw[sym] (0.95,1.25) -- (0.95,0.9) -- (0.25,0) -- (2.35,0) -- (1.65,0.9) -- (1.65,1.25);
    \fill[Ink] (1.3,0.75) circle (0.08);
  \or % 8 Distillation: flask, delivery tube, collecting tube
    \draw[sym] (0.5,2.6) -- (0.5,1.9) .. controls (-0.2,1.7) and (-0.2,0) .. (0.9,0) .. controls (2.0,0) and (2.0,1.7) .. (1.3,1.9) -- (1.3,2.6);
    \draw[sym] (0.5,2.6) -- (1.3,2.6);
    \fill[black!20] (0.15,0.8) .. controls (0.2,0.05) and (1.6,0.05) .. (1.65,0.8) -- cycle;
    \draw[Ink, line width=1.6pt] (1.3,2.3) -- (3.3,1.4) -- (3.3,0.9);
    \draw[sym] (3.0,1.5) -- (3.0,0.2) arc[start angle=180,end angle=360,radius=0.3] -- (3.6,1.5);
    \fill[Ink] (3.3,0.55) circle (0.08);
  \or % 9 The Great Green Mix Up: a mixture and a question
    \symbeaker
    \foreach \x/\y in {0.4/0.3,0.9/0.5,1.5/0.3,2.0/0.6,0.6/0.9,1.3/1.0,1.9/1.2,0.5/1.5,1.1/1.6}{\draw[Ink, line width=1pt, fill=white] (\x,\y) circle (0.15);}
    \foreach \x/\y in {0.7/0.35,1.2/0.35,1.75/0.35,2.2/0.3,1.6/0.75,0.9/1.25}{\fill[Ink] (\x-0.1,\y-0.1) rectangle (\x+0.1,\y+0.1);}
    \foreach \x/\y in {0.3/1.0,1.6/1.5,2.2/1.0}{\fill[black!40] (\x,\y) circle (0.08);}
    \node[font=\headingfont\bfseries\LARGE] at (3.3,2.2) {?};
  \fi
  \end{tikzpicture}}
\begin{document}
\setlength{\parindent}{0pt}
\setlength{\parskip}{4pt}

% ---- Cover: the unit's cover art, full bleed, in colour ------------------
% The art is 2:3 and A4 is taller-than-wide by less, so it fills the width and
% loses a little off the bottom; the baked-in title at the top stays whole.
\pdfbookmark[1]{Cover}{cover}
\thispagestyle{empty}
\stcover
\begin{tikzpicture}[remember picture, overlay]
  
  % Name box, top right, clear of the title.
  \node[anchor=north east, fill=white, draw=Ink, line width=0.6pt, rounded corners=1.5mm,
        inner xsep=2.6mm, inner ysep=2.2mm, text width=40mm, font=\headingfont\small]
    at ([xshift=-6mm, yshift=-6mm]current page.north east)
    {\textbf{Name}\par\vspace{5mm}\hrule\vspace{2.6mm}\textbf{Class}\par\vspace{5mm}\hrule};
  % What this book is: the cover art says Workbook, so name the booklet here.
  \node[anchor=south east, fill=white, draw=Ink, line width=0.6pt, rounded corners=1.5mm,
        inner xsep=3.2mm, inner ysep=2.4mm, font=\headingfont\bfseries\large, align=right]
    at ([xshift=-6mm, yshift=8mm]current page.south east)
    {Experiment Booklet\\[-1pt]{\normalfont\headingfont\small Experiments 1--9}};
\end{tikzpicture}
\mbox{}
\clearpage

% ---- Inside cover ---------------------------------------------------------
\thispagestyle{plain}
{\headingfont\bfseries\fontsize{24}{28}\selectfont Separation Technologies and Techniques}\par\vspace{1mm}
{\headingfont\Large Experiment Booklet \textbullet\ Year 7}\par\vspace{1mm}
{\small\itshape VC 2.0: VC2S8U05 \textbullet\ VC2S8U06 \textbullet\ VC2S8U07 \textbullet\ VC2S8I01--I08 \textbullet\ VC2S8H02}\par\vspace{6mm}
\renewcommand{\arraystretch}{2.2}
\begin{tabular}{@{}l p{130mm}@{}}
\textbf{Name} & \hrulefill \\ \textbf{Class} & \hrulefill \\ \textbf{Teacher} & \hrulefill \\ \textbf{Lab group} & \hrulefill \\
\end{tabular}
\renewcommand{\arraystretch}{1}

\begin{countrybox}{Acknowledgement of Country}
This booklet was created on the unceded Country of the Wurundjeri Woi-wurrung people of the Kulin Nation. We acknowledge Elders past, present and emerging. Aboriginal and Torres Strait Islander peoples have observed, tested and passed on knowledge of Country for tens of thousands of years: science in the way this booklet teaches it, by careful observation, measurement and sharing what was found.
\end{countrybox}

\begin{infobox}{How this booklet works}
This booklet holds every experiment in the unit. Bring it to every practical lesson. Each experiment opens with its key vocabulary, learning intentions and skills focus, then runs \emph{Technique} $\rightarrow$ \emph{Plan and safety} $\rightarrow$ \emph{Method} $\rightarrow$ \emph{Results} $\rightarrow$ \emph{Analysis and evaluation} $\rightarrow$ \emph{Skill check}. Each kind of box has its own frame and mark:\par\smallskip
\renewcommand{\arraystretch}{1.35}
\begin{tabular}{@{}p{32mm}l@{}}
\textbf{Technique} \tagicon{\icokey} & How to do it properly (thin frame) \\
\textbf{I do} \tagicon{\icodots{1}} & Your teacher works an example through --- you follow (solid frame) \\
\textbf{We do} \tagicon{\icodots{2}} & You try one together, with help (dashed frame) \\
\textbf{You do} \tagicon{\icodots{3}} & On your own: check your understanding (heavy frame) \\
\textbf{Going further} \tagicon{\icoup} & Above level: Year 8 extension (double frame) \\
\textbf{Skills focus} \tagicon{\icoskill} & The science skills the experiment builds \\
\end{tabular}

\smallskip
Every skill comes from the faculty \textbf{Science Skills} list (the toolkit on page \pageref{sec:toolkit}). A skill always carries the same code, for example \sk{PC4} or \sk{PA1}, so you can follow it from this unit into every other unit this year.
\end{infobox}

\clearpage
{\headingfont\bfseries\LARGE Contents}\par\vspace{1mm}\noindent\rule{\linewidth}{0.6pt}\par\vspace{2mm}
\begin{ctab}
\renewcommand{\arraystretch}{1.5}
\begin{tabular}{|L{2.8cm}|L{10.2cm}|C{1.6cm}|}
  \hline\thead{Section} & \thead{Title} & \thead{Page} \\ \hline
  Skills & \hyperlink{toolkit}{The Science Skills toolkit} & \hyperlink{toolkit}{\pageref{sec:toolkit}} \\ \hline
  Skills & \hyperlink{tracker}{My skills tracker} & \hyperlink{tracker}{\pageref{sec:tracker}} \\ \hline
  Unit & \hyperlink{glance}{Separation techniques at a glance} & \hyperlink{glance}{\pageref{sec:glance}} \\ \hline
  Experiment 1 & \hyperlink{exp.1}{Solubility (Lesson 3)} & \hyperlink{exp.1}{\pageref{sec:1}} \\ \hline
  Experiment 2 & \hyperlink{exp.2}{Surface Area and Dissolving (Lesson 6)} & \hyperlink{exp.2}{\pageref{sec:2}} \\ \hline
  Experiment 3 & \hyperlink{exp.3}{Fast and Slow Evaporation (Lesson 9)} & \hyperlink{exp.3}{\pageref{sec:3}} \\ \hline
  Experiment 4 & \hyperlink{exp.4}{Rates of Evaporation (Lesson 11)} & \hyperlink{exp.4}{\pageref{sec:4}} \\ \hline
  Experiment 5 & \hyperlink{exp.5}{Solubility and Temperature (Lesson 13)} & \hyperlink{exp.5}{\pageref{sec:5}} \\ \hline
  Experiment 6 & \hyperlink{exp.6}{Growing Copper Sulfate Crystals (Lesson 14)} & \hyperlink{exp.6}{\pageref{sec:6}} \\ \hline
  Experiment 7 & \hyperlink{exp.7}{Filtration (Lesson 15)} & \hyperlink{exp.7}{\pageref{sec:7}} \\ \hline
  Experiment 8 & \hyperlink{exp.8}{Distillation (Lesson 17)} & \hyperlink{exp.8}{\pageref{sec:8}} \\ \hline
  Experiment 9 & \hyperlink{exp.9}{The Great Green Mix Up (Lesson 18)} & \hyperlink{exp.9}{\pageref{sec:9}} \\ \hline
  Glossary & \hyperlink{glossary}{Key vocabulary, with the page each word first appears on} & \hyperlink{glossary}{\pageref{sec:glossary}} \\ \hline
\end{tabular}
\end{ctab}

% ══════════════════════════════════════════════════════════════════════════════
%  THE SCIENCE SKILLS TOOLKIT — faculty continuum, Years 7–8 column
% ══════════════════════════════════════════════════════════════════════════════
\clearpage
\hypertarget{toolkit}{}\pdfbookmark[1]{The Science Skills toolkit}{toolkitbm}
\markboth{The Science Skills toolkit}{}{\headingfont\bfseries\LARGE The Science Skills toolkit}\label{sec:toolkit}\par\vspace{1mm}\noindent\rule{\linewidth}{0.6pt}\par
These are the skills scientists use in every investigation. They come from the Science Inquiry strand of the Victorian Curriculum, grouped the way our faculty teaches them. Each experiment in this booklet tells you which skills it builds.

\newcommand{\tkhead}[2]{\multicolumn{3}{|l|}{\cellcolor{black!10}\rule{0pt}{4.6mm}\headingfont\bfseries #2} \\ \hline}
\newcommand{\tkrow}[1]{\skc{#1}\enspace{\small\bfseries\skname{#1}} & \skdesc{#1} & {\small\skcode{#1}} \\ \hline}
\renewcommand{\arraystretch}{1.35}
\begin{ctab}
\begin{longtable}{|W{5.0cm}|W{9.0cm}|Z{1.7cm}|}
  \hline\thead{Skill} & \thead{What it means} & \thead{Code} \\ \hline
  \endhead
  \tkhead{Teal}{Questioning and predicting}
  \tkrow{QP1}
  \tkrow{QP2}
  \tkhead{Blue}{Planning and conducting}
  \tkrow{PC1}
  \tkrow{PC2}
  \tkrow{PC3}
  \tkrow{PC4}
  \tkrow{PC5}
  \tkrow{PC6}
  \tkrow{PC7}
  \tkrow{PC8}
  \tkhead{GoldDark}{Processing, modelling and analysing}
  \tkrow{PA1}
  \tkrow{PA2}
  \tkrow{PA3}
  \tkrow{PA4}
  \tkrow{PA5}
  \tkrow{PA6}
  \tkrow{PA7}
  \tkrow{PA8}
  \tkhead{OrangeDark}{Evaluating}
  \tkrow{EV1}
  \tkrow{EV2}
  \tkrow{EV3}
  \tkrow{EV4}
  \tkrow{EV5}
  \tkhead{PurpleDark}{Communicating}
  \tkrow{CO1}
  \tkrow{CO2}
\end{longtable}
\end{ctab}
\renewcommand{\arraystretch}{1}

\begin{infobox}{Writing like a scientist \textbar\ CO1 \skname{CO1}}
Scientists write in the \textbf{third person}: say what was done, not who did it. Write \emph{``The water was heated for 10 minutes''}, not \emph{``I heated the water for 10 minutes''}. Keep sentences short, spell scientific words correctly, and always give a unit with a measurement.
\end{infobox}

% ══════════════════════════════════════════════════════════════════════════════
%  SKILLS TRACKER
% ══════════════════════════════════════════════════════════════════════════════
\clearpage
\hypertarget{tracker}{}\pdfbookmark[1]{My skills tracker}{trackerbm}
\markboth{My skills tracker}{}{\headingfont\bfseries\LARGE My skills tracker}\label{sec:tracker}\par\vspace{1mm}\noindent\rule{\linewidth}{0.6pt}\par
A dot \textbullet{} shows the experiments that build each skill. After each experiment, write how confident you are with each of its skills in its column: \textbf{1}~not yet, \textbf{2}~nearly, \textbf{3}~got it.

\newcommand{\dotc}{{\large\textbullet}}
\newcommand{\trow}[2]{\rule[-0.3cm]{0pt}{0.82cm}\skc{#1}\enspace{\footnotesize\skname{#1}} & #2 \\ \hline}
\setlength{\tabcolsep}{3pt}
\begin{ctab}
\begin{tabular}{|W{6.1cm}|*{9}{Z{0.95cm}|}}
  \hline\thead{Skill} & \thead{1} & \thead{2} & \thead{3} & \thead{4} & \thead{5} & \thead{6} & \thead{7} & \thead{8} & \thead{9} \\ \hline
  \trow{QP1}{\dotc & \dotc &  & \dotc &  &  &  &  & }
  \trow{QP2}{ &  &  & \dotc &  &  &  &  & \dotc}
  \trow{PC1}{ &  &  & \dotc &  &  &  &  & \dotc}
  \trow{PC2}{\dotc & \dotc &  & \dotc &  &  & \dotc &  & }
  \trow{PC3}{ &  & \dotc &  &  & \dotc &  & \dotc & \dotc}
  \trow{PC4}{ &  &  & \dotc & \dotc &  &  &  & }
  \trow{PC5}{ &  &  &  &  & \dotc &  &  & }
  \trow{PC6}{\dotc &  & \dotc &  & \dotc & \dotc & \dotc & \dotc & }
  \trow{PC7}{ &  &  &  &  &  & \dotc & \dotc & \dotc}
  \trow{PC8}{ & \dotc &  &  &  &  &  &  & }
  \trow{PA1}{\dotc & \dotc &  &  & \dotc &  &  &  & }
  \trow{PA2}{ &  &  & \dotc & \dotc &  &  &  & }
  \trow{PA3}{ & \dotc &  & \dotc &  &  &  &  & }
  \trow{PA5}{ & \dotc &  &  & \dotc &  &  &  & }
  \trow{PA6}{ &  & \dotc &  &  &  &  &  & }
  \trow{PA7}{ &  &  &  & \dotc &  &  &  & }
  \trow{EV1}{ & \dotc &  &  &  &  &  &  & }
  \trow{EV2}{ &  &  &  & \dotc &  &  & \dotc & }
  \trow{EV3}{ & \dotc &  & \dotc &  &  & \dotc &  & \dotc}
  \trow{EV5}{ &  & \dotc &  &  & \dotc &  &  & }
  \trow{CO1}{\dotc &  &  &  &  &  &  &  & \dotc}
  \trow{CO2}{ &  & \dotc &  &  & \dotc & \dotc & \dotc & }
\end{tabular}
\end{ctab}
\setlength{\tabcolsep}{6pt}
{\small\color{DarkGrey} Columns 1--9 are the experiments. The other skills are built in the experiment booklets for other units.}

% ══════════════════════════════════════════════════════════════════════════════
%  SEPARATION TECHNIQUES AT A GLANCE
% ══════════════════════════════════════════════════════════════════════════════
\clearpage
\hypertarget{glance}{}\pdfbookmark[1]{Separation techniques at a glance}{glancebm}
\markboth{Separation techniques at a glance}{}{\headingfont\bfseries\LARGE Separation techniques at a glance}\label{sec:glance}\par\vspace{1mm}\noindent\rule{\linewidth}{0.6pt}\par
Every separation uses a property that one part of a mixture has and the others do not. Use this page to plan a separation, and tick each technique when you have used it.
\renewcommand{\arraystretch}{1.45}
\begin{ctab}
\begin{tabular}{|L{2.9cm}|L{3.6cm}|L{4.4cm}|L{2.5cm}|C{1.3cm}|}
  \hline\thead{Technique} & \thead{Property it uses} & \thead{Separates} & \thead{Experiment} & \thead{Used} \\ \hline
  Sieving & Size of the pieces & Large solids from small solids & 9 & \tickbox \\ \hline
  Magnetic separation & Magnetic or not & Iron and steel from other solids & 9 & \tickbox \\ \hline
  Decanting & Insoluble solid sinks & A liquid from a settled solid & 7 & \tickbox \\ \hline
  Filtration & Insoluble; larger than the pores & An insoluble solid from a liquid & 7, 9 & \tickbox \\ \hline
  Evaporation & Solute does not evaporate & A dissolved solid from its solvent & 3, 4, 9 & \tickbox \\ \hline
  Crystallisation & Solubility falls as solvent leaves & Pure crystals from a solution & 3, 6 & \tickbox \\ \hline
  Distillation & Different boiling points & A pure liquid from a solution & 8 & \tickbox \\ \hline
  Chromatography & Different solubility in the solvent & The dyes in an ink & Lesson & \tickbox \\ \hline
\end{tabular}
\end{ctab}
\renewcommand{\arraystretch}{1}

\begin{infobox}{Pure substance or mixture?}
A \textbf{pure substance} contains only one kind of particle: an element (such as copper) or a compound (such as water). A \textbf{mixture} contains two or more substances that are not chemically joined, so it can be separated by physical methods like the ones above. A \textbf{solution} is a mixture in which a solute is dissolved in a solvent.
\end{infobox}

% ══════════════════════════════════════════════════════════════════════════════
%  EXPERIMENT 1 — SOLUBILITY                                     (Lesson 3, Term 2 Week 1)
% ══════════════════════════════════════════════════════════════════════════════
\experimentdiv{Solubility}{Testing which substances dissolve in water \textbullet{} Keeping a test fair \textbullet{} Recording observations in a table}{Lesson 3 \textbullet\ Term 2, Week 1 \textbullet\ Lab}\label{sec:1}
\lessoncodes{VC2S8U06 \textbullet{} VC2S8I01, I02, I04, I08}

\vocabhead
\begin{ctab}
\begin{tabular}{|SW{4.2cm}|W{11.6cm}|}
  \hline\hd & \thead{Word} & \thead{Definition (copy from the slide)} \\ \hline
  & \scaf{\textbf{Dissolve}\label{voc:dissolve}} & \scaf{When a solid mixes into a liquid so completely that it can no longer be seen.} \\ \hline
  & \textbf{Soluble}\label{voc:soluble} & \\ \hline
  & \textbf{Insoluble}\label{voc:insoluble} & \\ \hline
  & \textbf{Solution}\label{voc:solution} & \\ \hline
  & \textbf{Suspension}\label{voc:suspension} & \\ \hline
\end{tabular}
\end{ctab}

\begin{lipanel}
\inquiryli{Plan a fair test by keeping the amount of solid and water the same; record observations in a table.}
\understandli{Classify substances as soluble or insoluble in water.}
\sctitle
\begin{itemize}[leftmargin=1.4em, itemsep=0pt, topsep=1pt]
  \item I can predict which substances will dissolve and give a reason.
  \item I can name the variables I kept the same.
  \item I can classify each substance as soluble or insoluble from my observations.
\end{itemize}
\end{lipanel}

\begin{skillsfocus}
\focus{QP1}{Predict which substances will dissolve.}
\focus{PC2}{Keep the amounts of solid and water the same.}
\focus{PC6}{Observe each mixture closely.}
\focus{PA1}{Record observations in a table.}
\focus{CO1}{Write the conclusion in the third person.}
\end{skillsfocus}

% ─── TECHNIQUE ────────────────────────────────────────────────────────────────
\exppart{Technique spotlight}{Blue}

\begin{technique}{Telling soluble from insoluble}{PC6}
After stirring, hold the test tube up against a white card and look through it.
\begin{enumerate}[leftmargin=1.6em, itemsep=1pt, topsep=3pt]
  \item \textbf{Clear} (you can see through it, even if it is coloured) and no solid left: the substance is \textbf{soluble}. You have made a \textbf{solution}.
  \item \textbf{Cloudy}, and the solid settles to the bottom if you wait: the substance is \textbf{insoluble}. You have made a \textbf{suspension}.
  \item Clear, but some solid stays on the bottom: some of it dissolved. Record it as \textbf{partly soluble}.
\end{enumerate}
Clear does not mean colourless. Blue copper sulfate solution is clear; milky water is not.
\end{technique}

\begin{workedex}{Classifying from an observation}
\textbf{Observation:} \emph{After stirring, the liquid was cloudy brown. After one minute a layer of brown solid had settled on the bottom.}

\smallskip
\textbf{Classification:} insoluble. The liquid was cloudy and the solid settled out, so it is a suspension, not a solution.
\end{workedex}

\begin{wedo}{Classifying from an observation}
\emph{After stirring, the liquid was clear and pale yellow. No solid could be seen.} Classify the substance and give your reason.
\answerbox{3cm}
\end{wedo}

\begin{technique}{Keeping a test fair}{PC2}
In a fair test you change \textbf{one} thing (the independent variable), measure or observe \textbf{one} thing (the dependent variable), and keep \textbf{everything else} the same (the controlled variables). If you used more water for salt than for sand, you could not tell whether the salt dissolved because of what it is or because it had more water.
\end{technique}

\begin{youdobox}[Naming the variables]
\task Complete the table for today's experiment. \skr{PC2}
\begin{ctab}
\begin{tabular}{|SW{4.4cm}|W{10.2cm}|}
  \hline\hd & \thead{Variable} & \thead{In this experiment} \\ \hline
  & \scaf{Independent (what I change)} & \scaf{The substance being tested} \\ \hline
  & Dependent (what I observe) & \\ \hline
  & Controlled (what I keep the same) & \\ \hline
  & Controlled & \\ \hline
  & Controlled & \\ \hline
\end{tabular}
\end{ctab}
\end{youdobox}

% ─── PLAN AND SAFETY / METHOD ────────────────────────────────────────────────
\exppart{Plan, safety and method}{Teal}

\textbf{Aim:} To find out which common substances are soluble in water.

\task Predict whether each substance in the results table will dissolve in water. Write your prediction in the \emph{Prediction} column of Table 1.1, then explain one of your predictions here. \skr{QP1}
\answerbox{4cm}

\subsubsection*{Equipment}
8 test tubes in a rack \textbullet{} 10\,mL measuring cylinder \textbullet{} spatula \textbullet{} stirring rod \textbullet{} white card \textbullet{} small samples of salt, sugar, sand, chalk powder, flour, instant coffee, bicarbonate soda and cooking oil \textbullet{} wash bottle of water

\begin{center}
\begin{tikzpicture}[scale=0.8]
  \begin{scope}[shift={(0,0)}]\eqTestTube\fill[black!8] (0.02,0.3) arc[start angle=180,end angle=360,radius=0.28] -- (0.58,1.6) -- (0.02,1.6) -- cycle;\end{scope}
  \begin{scope}[shift={(2.2,0)}]\eqTestTube\fill[black!30] (0.02,0.3) arc[start angle=180,end angle=360,radius=0.28] -- (0.58,1.6) -- (0.02,1.6) -- cycle;\end{scope}
  \begin{scope}[shift={(4.4,0)}]\eqTestTube\fill[black!70] (0.02,0.3) arc[start angle=180,end angle=360,radius=0.28] -- (0.58,0.5) -- (0.02,0.5) -- cycle;\end{scope}
  \foreach \x in {0,2.2,4.4}{\draw[Ink, line width=0.8pt] (\x,1.6) -- (\x+0.6,1.6);}
  \node[eqtext, align=center] at (0.3,-0.45) {clear:\\soluble};
  \node[eqtext, align=center] at (2.5,-0.45) {cloudy:\\suspension};
  \node[eqtext, align=center] at (4.7,-0.45) {solid settled:\\insoluble};
\end{tikzpicture}
\expfigcap{What to look for after stirring.}
\end{center}

\begin{infobox}{Safety}
Wear safety glasses. Never taste anything in the lab, even substances from the kitchen. Wipe up spills of powder or oil straight away so that nobody slips.
\end{infobox}

\subsubsection*{Method}
\begin{enumerate}[itemsep=2pt]
  \item Label eight test tubes with the names of the substances.
  \item Measure 10\,mL of water into each test tube.
  \item Add one level spatula of the first substance to its test tube.
  \item Stir for 30 seconds, then hold the test tube against the white card and observe it.
  \item Wait one minute and observe again. Record both observations in Table 1.1.
  \item Repeat steps 3 to 5 for each substance, using a clean stirring rod each time.
\end{enumerate}

% ─── RESULTS ─────────────────────────────────────────────────────────────────
\exppart{Results}{Amber}

\subsubsection*{Table 1.1: Do these substances dissolve in water? \skr{PA1}}
\begin{ctab}
\begin{tabular}{|SW{2.6cm}|W{2.2cm}|W{3.8cm}|W{3.2cm}|W{2.6cm}|}
  \hline\hd & \thead{Substance} & \thead{Prediction} & \thead{After stirring} & \thead{After 1 minute} & \thead{Soluble or insoluble?} \\ \hline
  & \scaf{Salt} & \scaf{Will dissolve} & \scaf{Clear, colourless; no solid seen} & \scaf{No change} & \scaf{Soluble} \\ \hline
  & Sugar & & & & \\ \hline
  & Sand & & & & \\ \hline
  & Chalk powder & & & & \\ \hline
  & Flour & & & & \\ \hline
  & Instant coffee & & & & \\ \hline
  & Bicarbonate soda & & & & \\ \hline
  & Cooking oil & & & & \\ \hline
\end{tabular}
\end{ctab}

% ─── ANALYSIS ────────────────────────────────────────────────────────────────
\exppart{Analysis and evaluation}{Coral}

\task Sort the substances you tested into two lists. \skr{PA1}
\begin{ctab}
\begin{tabular}{|W{7.4cm}|W{7.4cm}|}
  \hline\thead{Soluble in water} & \thead{Insoluble in water} \\ \hline
  \rule{0pt}{3.2cm} & \\ \hline
\end{tabular}
\end{ctab}

\needspace{10\baselineskip}
\task How many of your predictions were correct? Choose one that was wrong and suggest why you predicted it wrongly.
\answerbox{4cm}

\needspace{10\baselineskip}
\task Instant coffee made a dark liquid. A student said, ``It's dark, so it can't have dissolved.'' Explain why the student is wrong.
\answerbox{4cm}

\needspace{10\baselineskip}
\task Explain why it was important to use the same volume of water and the same amount of each solid. \skr{PC2}
\answerbox{4cm}

\needspace{12\baselineskip}
\task Write a conclusion for this experiment. Answer the aim, give two examples, and write in the third person (no \emph{I} or \emph{we}). \skr{CO1}
\answerbox{5cm}

\begin{yr8box}
\begin{keepext}
\exttask Oil did not dissolve in water, but it does dissolve in other liquids, such as eucalyptus oil. Plan a fair test to find out which of three liquids dissolves the most substances. Name your independent, dependent and controlled variables.
\answerbox{6cm}
\end{keepext}
\end{yr8box}

% ─── SKILL CHECK ─────────────────────────────────────────────────────────────
\exppart{Skill check}{NavyBlue}

\begin{skillcheck}
\scrow{QP1}{I can predict which substances will dissolve and give a reason.}
\scrow{PC2}{I can name the variables I kept the same.}
\scrow{PC6}{I can tell a solution from a suspension by looking.}
\scrow{PA1}{I recorded every observation in my table.}
\scrow{CO1}{I can write a conclusion in the third person.}
\end{skillcheck}

\task Which skill will you focus on in Experiment 2? Explain what you will do differently.
\answerbox{5cm}
% ══════════════════════════════════════════════════════════════════════════════
%  EXPERIMENT 2 — SURFACE AREA AND DISSOLVING                    (Lesson 6, Term 2 Week 2)
% ══════════════════════════════════════════════════════════════════════════════
\experimentdiv{Surface Area and Dissolving}{Timing how fast sugar dissolves \textbullet{} Repeat trials and averages \textbullet{} Testing a hypothesis}{Lesson 6 \textbullet\ Term 2, Week 2 \textbullet\ Lab}\label{sec:2}
\lessoncodes{VC2S8U05, U06 \textbullet{} VC2S8I01--I06}

\vocabhead
\begin{ctab}
\begin{tabular}{|SW{4.2cm}|W{11.6cm}|}
  \hline\hd & \thead{Word} & \thead{Definition (copy from the slide)} \\ \hline
  & \scaf{\textbf{Solute}\label{voc:solute}} & \scaf{The substance that dissolves, such as sugar.} \\ \hline
  & \textbf{Solvent}\label{voc:solvent} & \\ \hline
  & \textbf{Surface area}\label{voc:surface-area} & \\ \hline
  & \textbf{Rate}\label{voc:rate} & \\ \hline
  & \textbf{Average}\label{voc:average} & \\ \hline
\end{tabular}
\end{ctab}

\begin{lipanel}
\inquiryli{Write a hypothesis, time repeat trials and calculate an average.}
\understandli{Explain, using the particle model, why smaller pieces of solute dissolve faster.}
\sctitle
\begin{itemize}[leftmargin=1.4em, itemsep=0pt, topsep=1pt]
  \item I can write an \emph{If \dots\ then \dots\ because \dots} hypothesis.
  \item I can time three trials and calculate their average.
  \item I can describe the trend in my results and say whether it supports my hypothesis.
\end{itemize}
\end{lipanel}

\begin{skillsfocus}
\focus{QP1}{Write a testable hypothesis.}
\focus{PC2}{Name the independent, dependent and controlled variables.}
\focus{PC8}{Time each trial with a digital stopwatch.}
\focus{PA1}{Record times in a results table.}
\focus{PA3}{Calculate the average of three trials.}
\focus{PA5}{Describe the trend.}
\focus{EV1}{Compare the results with the hypothesis.}
\focus{EV3}{Find sources of error.}
\end{skillsfocus}

% ─── TECHNIQUE ────────────────────────────────────────────────────────────────
\exppart{Technique spotlight}{Blue}

\begin{technique}{Writing a hypothesis}{QP1}
A \textbf{hypothesis} is a prediction you can test, with a reason. Write it in three parts:
\begin{center}\emph{If} [I change the independent variable this way], \emph{then} [the dependent variable will do this], \emph{because} [the science reason].\end{center}
\end{technique}

\begin{workedex}{A hypothesis about temperature}
\emph{If the water is hotter, then sugar will dissolve in less time, because the water particles move faster and hit the sugar more often.}
\end{workedex}

\begin{wedo}{A hypothesis about stirring}
Write a hypothesis for an experiment that tests whether stirring changes how fast salt dissolves.
\answerbox{3.5cm}
\end{wedo}

\begin{technique}{Repeat trials and averages}{PA3}
One timing can be wrong by chance: you might start the stopwatch late. Do each test \textbf{three times} and calculate the \textbf{average}:
\begin{center}average $=$ (trial 1 $+$ trial 2 $+$ trial 3) $\div$ 3\end{center}
If one trial is very different from the other two, it is an \textbf{outlier}. Circle it, leave it out, and average the other two.
\end{technique}

\begin{workedex}{Calculating an average}
Times: 42\,s, 45\,s, 39\,s.\qquad Average $= (42 + 45 + 39) \div 3 = 126 \div 3 = 42$\,s.
\end{workedex}

\begin{youdobox}[Averages]
\task Calculate the average of each set of times. Show your working. \skr{PA3}
\pt{a}55\,s, 61\,s, 58\,s
\workbox{2cm}
\pt{b}20\,s, 23\,s, 47\,s (look for an outlier first)
\workbox{2.5cm}
\end{youdobox}

% ─── PLAN AND SAFETY / METHOD ────────────────────────────────────────────────
\exppart{Plan, safety and method}{Teal}

\textbf{Aim:} To find out how the size of the sugar pieces affects the time sugar takes to dissolve.

\task Name the variables in this experiment. \skr{PC2}
\begin{ctab}
\begin{tabular}{|SW{4.4cm}|W{10.8cm}|}
  \hline\hd & \thead{Variable} & \thead{In this experiment} \\ \hline
  & \scaf{Independent} & \scaf{Size of the sugar pieces (cube, granulated, caster)} \\ \hline
  & Dependent & \\ \hline
  & Controlled & \\ \hline
  & Controlled & \\ \hline
\end{tabular}
\end{ctab}

\needspace{8\baselineskip}
\task Write your hypothesis. \skr{QP1}
\answerbox{2.4cm}

\subsubsection*{Equipment}
3 $\times$ 250\,mL beakers \textbullet{} 100\,mL measuring cylinder \textbullet{} electronic balance \textbullet{} stirring rod \textbullet{} digital stopwatch (or phone) \textbullet{} sugar cubes, granulated sugar and caster sugar \textbullet{} room-temperature water

\begin{infobox}{Safety}
Wear safety glasses. Do not eat the sugar: lab equipment is not clean enough for food. Keep the phone away from the water.
\end{infobox}

\subsubsection*{Method}
\begin{enumerate}[itemsep=2pt]
  \item Measure 100\,mL of water into a beaker.
  \item Weigh out 4\,g of sugar: one cube, or 4\,g of granulated or caster sugar.
  \item Tip the sugar into the water and start the stopwatch at the same moment.
  \item Stir at a steady rate, about one stir per second. Stop the stopwatch when no sugar can be seen.
  \item Record the time in Table 2.1, then rinse the beaker.
  \item Repeat for the other two sizes, then do the whole experiment twice more.
\end{enumerate}

% ─── RESULTS ─────────────────────────────────────────────────────────────────
\exppart{Results}{Amber}

\subsubsection*{Table 2.1: Time for 4\,g of sugar to dissolve in 100\,mL of water \skr{PA1}}
\begin{ctab}
\begin{tabular}{|SW{3.6cm}|C{2.3cm}|C{2.3cm}|C{2.3cm}|C{3.0cm}|}
  \hline\hd & \thead{Size of pieces} & \thead{Trial 1 (s)} & \thead{Trial 2 (s)} & \thead{Trial 3 (s)} & \thead{Average (s)} \\ \hline
  & \scaf{Sugar cube} & \scaf{150} & \scaf{162} & \scaf{141} & \scaf{151} \\ \hline
  & Sugar cube & & & & \\ \hline
  & Granulated sugar & & & & \\ \hline
  & Caster sugar & & & & \\ \hline
\end{tabular}
\end{ctab}
{\small\color{DarkGrey}The red row is an example. Record your own results in the rows below it.}

\task Show your working for one of the averages. \skr{PA3}
\workbox{3cm}

\needspace{16\baselineskip}
\task Draw a column graph of the average times. Put the size of pieces on the horizontal axis. Label both axes and give the graph a title. \skr{PA2}
\graphgrid{8cm}

% ─── ANALYSIS ────────────────────────────────────────────────────────────────
\exppart{Analysis and evaluation}{Coral}

\task Describe the trend in your results. Use numbers from your table. \skr{PA5}
\answerbox{4cm}

\needspace{10\baselineskip}
\task Do your results support your hypothesis? Explain. \skr{EV1}
\answerbox{4cm}

\needspace{12\baselineskip}
\task Use the particle model to explain why caster sugar dissolves faster than a sugar cube. Use the words \emph{surface area}, \emph{solvent particles} and \emph{solute particles}.
\answerbox{5cm}

\needspace{12\baselineskip}
\task Give two sources of error in this experiment, and say how each could be reduced. \skr{EV3}
\begin{ctab}
\begin{tabular}{|SW{7.2cm}|W{7.8cm}|}
  \hline\hd & \thead{Source of error} & \thead{How to reduce it} \\ \hline
  & \scaf{Deciding when the last sugar has gone} & \scaf{Have the same person judge every trial, looking against a white card} \\ \hline
  & & \\ \hline
  & & \\ \hline
\end{tabular}
\end{ctab}

\begin{yr8box}
\begin{keepext}
\exttask A sugar cube is 1\,cm on each side. Calculate its surface area. If it is cut into eight cubes 0.5\,cm on each side, calculate the total surface area. What does this show about cutting a solid into smaller pieces?
\workbox{5cm}
\end{keepext}
\end{yr8box}

% ─── SKILL CHECK ─────────────────────────────────────────────────────────────
\exppart{Skill check}{NavyBlue}

\begin{skillcheck}
\scrow{QP1}{I can write an If \dots\ then \dots\ because \dots\ hypothesis.}
\scrow{PC2}{I can name the independent, dependent and controlled variables.}
\scrow{PC8}{I timed each trial with a digital stopwatch.}
\scrow{PA1}{I recorded my times in a table with units in the headings.}
\scrow{PA3}{I can calculate an average and spot an outlier.}
\scrow{PA5}{I can describe a trend using numbers.}
\scrow{EV1}{I can say whether my results support my hypothesis.}
\scrow{EV3}{I can name sources of error and how to reduce them.}
\end{skillcheck}

\task Which skill will you focus on in Experiment 3? Explain what you will do differently.
\answerbox{5cm}
% ══════════════════════════════════════════════════════════════════════════════
%  EXPERIMENT 3 — FAST AND SLOW EVAPORATION                      (Lesson 9, Term 2 Week 3)
% ══════════════════════════════════════════════════════════════════════════════
\experimentdiv{Fast and Slow Evaporation}{Evaporating to dryness safely \textbullet{} Observing crystals with a hand lens \textbullet{} Linking cause and effect}{Lesson 9 \textbullet\ Term 2, Week 3 \textbullet\ Lab}\label{sec:3}
\lessoncodes{VC2S8U05, U07 \textbullet{} VC2S8I02, I05, I07, I08}

\vocabhead
\begin{ctab}
\begin{tabular}{|SW{4.2cm}|W{11.6cm}|}
  \hline\hd & \thead{Word} & \thead{Definition (copy from the slide)} \\ \hline
  & \scaf{\textbf{Evaporation}\label{voc:evaporation}} & \scaf{A liquid changing to a gas at its surface, below its boiling point.} \\ \hline
  & \textbf{Crystal}\label{voc:crystal} & \\ \hline
  & \textbf{Crystallisation}\label{voc:crystallisation} & \\ \hline
  & \textbf{Residue}\label{voc:residue} & \\ \hline
  & \textbf{Cause and effect}\label{voc:cause-and-effect} & \\ \hline
\end{tabular}
\end{ctab}

\begin{lipanel}
\inquiryli{Heat a solution safely and draw what you observe with a hand lens.}
\understandli{Explain how the rate of evaporation affects the size of the crystals that form.}
\sctitle
\begin{itemize}[leftmargin=1.4em, itemsep=0pt, topsep=1pt]
  \item I can heat a solution safely and stop at the right time.
  \item I can draw a labelled diagram of the set-up and of the crystals.
  \item I can explain why slow evaporation makes larger crystals.
\end{itemize}
\end{lipanel}

\begin{skillsfocus}
\focus{PC3}{Heat an evaporating basin safely.}
\focus{PC6}{Observe the crystals with a hand lens.}
\focus{CO2}{Draw the set-up and the crystals.}
\focus{PA6}{Link the rate of evaporation to crystal size.}
\focus{EV5}{Suggest a question for a further investigation.}
\end{skillsfocus}

% ─── TECHNIQUE ────────────────────────────────────────────────────────────────
\exppart{Technique spotlight}{Blue}

\begin{technique}{Evaporating a solution with a Bunsen burner}{PC3}
\begin{enumerate}[leftmargin=1.6em, itemsep=1pt, topsep=3pt]
  \item Stand the evaporating basin on a gauze mat on a tripod, with a heatproof mat under the Bunsen.
  \item Use a \textbf{blue} flame, about half open. Stand up, wear safety glasses, and tie back long hair.
  \item Heat gently. When the first crystals appear at the edge, \textbf{turn the Bunsen off}. The hot basin finishes the job; heating to dryness makes the salt spit out.
  \item Leave the basin to cool on the gauze. Move it only with tongs once it has cooled.
\end{enumerate}
\end{technique}

\begin{workedex}{Drawing a scientific diagram}
A scientific diagram is a 2D side view drawn in pencil with a ruler. Every label sits to the side of the drawing, joined to its part by a straight line that does not cross another line.

\smallskip
\begin{center}
\begin{tikzpicture}[scale=0.62]
  \begin{scope}[shift={(0,0)}]\eqBunsen\end{scope}
  \begin{scope}[shift={(-0.4,2.4)}]\eqTripod\end{scope}
  \draw[eq] (-0.6,4.52) -- (2.4,4.52);
  \begin{scope}[shift={(-0.4,4.52)}]\eqBasin\end{scope}
  \fill[black!15] (-0.3,4.9) .. controls (0.2,4.6) and (1.6,4.6) .. (2.1,4.9) -- cycle;
  \eqlab{(2.0,5.3)}{(3.6,5.8)}{evaporating basin}
  \eqlab{(2.3,4.52)}{(3.6,4.7)}{gauze mat}
  \eqlab{(2.0,3.5)}{(3.6,3.5)}{tripod}
  \eqlab{(1.0,1.6)}{(3.6,1.8)}{Bunsen burner}
\end{tikzpicture}
\end{center}
\end{workedex}

\begin{wedo}{Spot the mistakes}
A student drew the set-up in pen, shaded the flame, and wrote the labels on the drawing. List what they should change.
\answerbox{3.5cm}
\end{wedo}

% ─── PLAN AND SAFETY / METHOD ────────────────────────────────────────────────
\exppart{Plan, safety and method}{Teal}

\textbf{Aim:} To compare the crystals that form when a salt solution evaporates quickly and slowly.

\task Predict which method will make the larger crystals. Give a reason. \skr{QP1}
\answerbox{4cm}

\subsubsection*{Equipment}
Saturated salt solution \textbullet{} 10\,mL measuring cylinder \textbullet{} evaporating basin \textbullet{} watch glass \textbullet{} Bunsen burner, tripod, gauze mat and heatproof mat \textbullet{} tongs \textbullet{} hand lens \textbullet{} safety glasses

\begin{infobox}{Safety}
Hot glass and hot basins look exactly like cold ones: touch them only with tongs. Turn the Bunsen off as soon as crystals appear, so that hot salt does not spit. Keep the gas off when you are not using it.
\end{infobox}

\subsubsection*{Method}
\textbf{Fast evaporation}
\begin{enumerate}[itemsep=2pt]
  \item Pour 10\,mL of salt solution into the evaporating basin and set up the Bunsen as in the technique box.
  \item Heat gently until crystals start to form at the edge. Turn off the gas and leave the basin to cool.
\end{enumerate}
\textbf{Slow evaporation}
\begin{enumerate}[itemsep=2pt, resume]
  \item Pour 5\,mL of salt solution onto a watch glass and label it with your names.
  \item Leave it on the windowsill until next lesson.
\end{enumerate}
\textbf{Both}
\begin{enumerate}[itemsep=2pt, resume]
  \item Look at each set of crystals through the hand lens. Record your observations and draw them.
\end{enumerate}

% ─── RESULTS ─────────────────────────────────────────────────────────────────
\exppart{Results}{Amber}

\task Draw a scientific diagram of the fast-evaporation set-up. \skr{CO2}
\drawbox{7cm}{}

\subsubsection*{Table 3.1: Crystals from fast and slow evaporation \skr{PC6}}
\begin{ctab}
\begin{tabular}{|SW{3.4cm}|W{3.4cm}|W{3.6cm}|W{4.6cm}|}
  \hline\hd & \thead{Method} & \thead{Time taken} & \thead{Size of crystals} & \thead{Shape and appearance} \\ \hline
  & \scaf{Example: sugar, fast} & \scaf{5 minutes} & \scaf{Too small to see one by one} & \scaf{White crust, like powder} \\ \hline
  & Fast (Bunsen) & & & \\ \hline
  & Slow (windowsill) & & & \\ \hline
\end{tabular}
\end{ctab}

\needspace{14\baselineskip}
\task Draw what you saw through the hand lens. \skr{CO2}

\noindent
\begin{minipage}[t]{0.48\linewidth}\drawbox{6cm}{}\nameline{Fast evaporation}\end{minipage}\hfill
\begin{minipage}[t]{0.48\linewidth}\drawbox{6cm}{}\nameline{Slow evaporation}\end{minipage}

% ─── ANALYSIS ────────────────────────────────────────────────────────────────
\exppart{Analysis and evaluation}{Coral}

\task Was your prediction correct? Use your observations to answer.
\answerbox{3.5cm}

\needspace{12\baselineskip}
\task Complete the sentence: \emph{The faster a solution evaporates, the \dots\ the crystals.} Then use the particle model to explain why. \skr{PA6}
\answerbox{5cm}

\needspace{10\baselineskip}
\task The salt crystals were white, but there was no sign of water in the basin at the end. Where did the water go? Name the change of state.
\answerbox{3.5cm}

\needspace{10\baselineskip}
\task Salt farms in South Australia make salt by letting seawater evaporate in large shallow ponds. Explain why the ponds are wide and shallow rather than deep.
\answerbox{4cm}

\needspace{10\baselineskip}
\task Write one question about crystals that you could test in another experiment. \skr{EV5}
\answerbox{3cm}

\begin{yr8box}
\begin{keepext}
\exttask A student says, ``The crystals were bigger on the windowsill because the windowsill was sunny.'' Is this a cause, or only something that happened at the same time? Plan how you could test it.
\answerbox{5cm}
\end{keepext}
\end{yr8box}

% ─── SKILL CHECK ─────────────────────────────────────────────────────────────
\exppart{Skill check}{NavyBlue}

\begin{skillcheck}
\scrow{PC3}{I heated the solution safely and stopped at the right time.}
\scrow{PC6}{I observed the crystals closely with a hand lens.}
\scrow{CO2}{I drew a scientific diagram of the set-up and the crystals.}
\scrow{PA6}{I can link the rate of evaporation to the size of the crystals.}
\scrow{EV5}{I can suggest a question for a further investigation.}
\end{skillcheck}

\task Which skill will you focus on in Experiment 4? Explain what you will do differently.
\answerbox{5cm}
% ══════════════════════════════════════════════════════════════════════════════
%  EXPERIMENT 4 — RATES OF EVAPORATION                           (Lesson 11, Term 2 Week 4)
% ══════════════════════════════════════════════════════════════════════════════
\experimentdiv{Rates of Evaporation}{Writing an aim \textbullet{} Planning your own method \textbullet{} Measuring mass loss over time}{Lesson 11 \textbullet\ Term 2, Week 4 \textbullet\ Lab}\label{sec:4}
\lessoncodes{VC2S8U05 \textbullet{} VC2S8I01--I04, I06}

\vocabhead
\begin{ctab}
\begin{tabular}{|SW{4.2cm}|W{11.6cm}|}
  \hline\hd & \thead{Word} & \thead{Definition (copy from the slide)} \\ \hline
  & \scaf{\textbf{Aim}\label{voc:aim}} & \scaf{A sentence saying what the experiment is trying to find out.} \\ \hline
  & \textbf{Rate of evaporation}\label{voc:rate-of-evaporation} & \\ \hline
  & \textbf{Method}\label{voc:method} & \\ \hline
  & \textbf{Repeatable}\label{voc:repeatable} & \\ \hline
  & \textbf{Mass lost}\label{voc:mass-lost} & \\ \hline
\end{tabular}
\end{ctab}

\begin{lipanel}
\inquiryli{Write an aim and a hypothesis, then plan and follow a repeatable method.}
\understandli{Explain how one factor affects the rate of evaporation.}
\sctitle
\begin{itemize}[leftmargin=1.4em, itemsep=0pt, topsep=1pt]
  \item I can write an aim and a hypothesis for my own investigation.
  \item I can write numbered method steps that someone else could follow.
  \item I can calculate mass lost and use it to compare rates.
\end{itemize}
\end{lipanel}

\begin{skillsfocus}
\focus{QP2}{Write an aim for your investigation.}
\focus{QP1}{Write a hypothesis with a reason.}
\focus{PC1}{Plan a method someone else could repeat.}
\focus{PC2}{Choose the variable to change and those to keep the same.}
\focus{PC4}{Weigh to 0.1\,g on the same balance each time.}
\focus{PA2}{Graph the mass lost.}
\focus{PA3}{Calculate the mass lost.}
\focus{EV3}{Find sources of error.}
\end{skillsfocus}

% ─── TECHNIQUE ────────────────────────────────────────────────────────────────
\exppart{Technique spotlight}{Blue}

\begin{technique}{Writing an aim}{QP2}
An aim names the independent and dependent variables and how they are linked:
\begin{center}\emph{To find out how} [the independent variable] \emph{affects} [the dependent variable].\end{center}
\end{technique}

\begin{workedex}{An aim}
\emph{To find out how the temperature of water affects the time sugar takes to dissolve.}
\end{workedex}

\begin{wedo}{An aim}
Write an aim for an experiment that tests whether a fan changes how quickly a wet towel dries.
\answerbox{3cm}
\end{wedo}

\begin{technique}{Writing a repeatable method}{PC1}
A method is \textbf{repeatable} when another group could follow it and get similar results. Number each step, start each step with a verb (\emph{Measure}, \emph{Place}, \emph{Record}), and give every amount with its unit.
\end{technique}

\begin{workedex}{Vague step and precise step}
\textbf{Vague:} \emph{Put some water in the dish.}\qquad
\textbf{Precise:} \emph{Measure 20\,mL of water and pour it into the Petri dish.}
\end{workedex}

\begin{youdobox}[Precise steps]
\task Rewrite each step so that it is precise. \skr{PC1}
\pt{a}Weigh it.
\answerbox{2cm}
\pt{b}Leave it somewhere for a while.
\answerbox{2cm}
\end{youdobox}

% ─── PLAN ────────────────────────────────────────────────────────────────────
\exppart{Plan your investigation}{Teal}

Choose \textbf{one} factor to test. Tick it.

\medskip
\begin{tabular}{@{}l l@{}}
\tickbox\enspace Surface area (Petri dish, beaker, test tube) & \tickbox\enspace Air movement (fan on, fan off) \\[3pt]
\tickbox\enspace Temperature (warm spot, cool spot) & \tickbox\enspace Covered or uncovered \\
\end{tabular}

\task Write your aim. \skr{QP2}
\answerbox{3cm}

\needspace{8\baselineskip}
\task Write your hypothesis. \skr{QP1}
\answerbox{3.5cm}

\needspace{10\baselineskip}
\task Name your variables. \skr{PC2}
\begin{ctab}
\begin{tabular}{|SW{4.4cm}|W{10.8cm}|}
  \hline\hd & \thead{Variable} & \thead{In my investigation} \\ \hline
  & \scaf{Dependent} & \scaf{Mass of water lost (g)} \\ \hline
  & Independent & \\ \hline
  & Controlled & \\ \hline
  & Controlled & \\ \hline
  & Controlled & \\ \hline
\end{tabular}
\end{ctab}

\subsubsection*{Equipment}
Containers for your factor \textbullet{} 50\,mL measuring cylinder \textbullet{} electronic balance (0.1\,g) \textbullet{} water \textbullet{} marker and masking tape \textbullet{} anything else your factor needs (fan, plastic wrap, a warm and a cool place)

\needspace{14\baselineskip}
\task Write your method as numbered steps. Include how you will weigh the containers and how often. \skr{PC1}
\answerbox{9cm}

\begin{infobox}{Safety}
Keep the water well away from the balance's power cable and from any fan. Wipe up spills straight away. Have your teacher check your method before you start.\par\smallskip
Teacher's initials: \rule{2.5cm}{0.4pt}
\end{infobox}

% ─── RESULTS ─────────────────────────────────────────────────────────────────
\exppart{Results}{Amber}

\subsubsection*{Table 4.1: Mass of each container and its water \skr{PC4}}
\begin{ctab}
\begin{tabular}{|SW{2.8cm}|C{2.2cm}|C{2.2cm}|C{2.2cm}|C{2.2cm}|C{2.2cm}|}
  \hline\hd & \thead{Container or condition} & \thead{Start (g)} & \thead{Lesson 2 (g)} & \thead{Lesson 3 (g)} & \thead{Lesson 4 (g)} & \thead{Mass lost (g)} \\ \hline
  & \scaf{Petri dish} & \scaf{62.4} & \scaf{58.1} & \scaf{53.7} & \scaf{49.5} & \scaf{12.9} \\ \hline
  & & & & & & \\ \hline
  & & & & & & \\ \hline
  & & & & & & \\ \hline
\end{tabular}
\end{ctab}

\task Show how you calculated one mass lost. \skr{PA3}
\workbox{2.5cm}

\needspace{18\baselineskip}
\task Draw a column graph of the mass lost by each container or condition. Label both axes and give the graph a title. \skr{PA2}
\graphgrid{9cm}

% ─── ANALYSIS ────────────────────────────────────────────────────────────────
\exppart{Analysis and evaluation}{Coral}

\task Which container or condition had the fastest rate of evaporation? How do you know? \skr{PA3}
\answerbox{3.5cm}

\needspace{10\baselineskip}
\task Do your results support your hypothesis? Use numbers from Table 4.1. \skr{EV1}
\answerbox{4cm}

\needspace{12\baselineskip}
\task Explain your results using the particle model.
\answerbox{5cm}

\needspace{12\baselineskip}
\task Give two sources of error and how each could be reduced. \skr{EV3}
\begin{ctab}
\begin{tabular}{|SW{7.2cm}|W{7.8cm}|}
  \hline\hd & \thead{Source of error} & \thead{How to reduce it} \\ \hline
  & \scaf{Water spilled when moving the container} & \scaf{Weigh the container where it stands, or carry it on a tray} \\ \hline
  & & \\ \hline
  & & \\ \hline
\end{tabular}
\end{ctab}

\needspace{10\baselineskip}
\task Was your method repeatable? Name one step another group found hard to follow, and rewrite it. \skr{PC1}
\answerbox{4cm}

\begin{yr8box}
\begin{keepext}
\exttask Rate means \emph{how much per unit of time}. Calculate the rate of evaporation of your fastest container in grams per day, and use it to predict how many days it would take for all its water to evaporate.
\workbox{5cm}
\end{keepext}
\end{yr8box}

% ─── SKILL CHECK ─────────────────────────────────────────────────────────────
\exppart{Skill check}{NavyBlue}

\begin{skillcheck}
\scrow{QP2}{I can write an aim that names both variables.}
\scrow{QP1}{I can write a hypothesis with a reason.}
\scrow{PC1}{I wrote a method another group could repeat.}
\scrow{PC2}{I changed one variable and kept the others the same.}
\scrow{PC4}{I weighed to 0.1\,g on the same balance each time.}
\scrow{PA2}{I drew a column graph of my results.}
\scrow{PA3}{I calculated mass lost correctly.}
\scrow{EV3}{I can name sources of error and how to reduce them.}
\end{skillcheck}

\task Which skill will you focus on in Experiment 5? Explain what you will do differently.
\answerbox{5cm}
% ══════════════════════════════════════════════════════════════════════════════
%  EXPERIMENT 5 — SOLUBILITY AND TEMPERATURE                     (Lesson 13, Term 2 Week 5)
% ══════════════════════════════════════════════════════════════════════════════
\experimentdiv{Solubility and Temperature}{Making a saturated solution \textbullet{} Pooling class data \textbullet{} Drawing a solubility curve}{Lesson 13 \textbullet\ Term 2, Week 5 \textbullet\ Lab}\label{sec:5}
\lessoncodes{VC2S8U05, U06 \textbullet{} VC2S8I02, I04, I05, I06}

\vocabhead
\begin{ctab}
\begin{tabular}{|SW{4.2cm}|W{11.6cm}|}
  \hline\hd & \thead{Word} & \thead{Definition (copy from the slide)} \\ \hline
  & \scaf{\textbf{Saturated}\label{voc:saturated}} & \scaf{A solution that holds as much solute as it can at that temperature; any more stays undissolved.} \\ \hline
  & \textbf{Solubility}\label{voc:solubility} & \\ \hline
  & \textbf{Solubility curve}\label{voc:solubility-curve} & \\ \hline
  & \textbf{Anomaly}\label{voc:anomaly} & \\ \hline
  & \textbf{Line of best fit}\label{voc:line-of-best-fit} & \\ \hline
\end{tabular}
\end{ctab}

\begin{lipanel}
\inquiryli{Measure precisely, pool class data and draw a line graph with a line of best fit.}
\understandli{Describe how the solubility of a solid changes with temperature.}
\sctitle
\begin{itemize}[leftmargin=1.4em, itemsep=0pt, topsep=1pt]
  \item I can tell when a solution is saturated.
  \item I can plot class data and draw a smooth line of best fit.
  \item I can spot an anomaly and suggest a cause.
\end{itemize}
\end{lipanel}

\begin{skillsfocus}
\focus{PC4}{Measure mass and temperature precisely.}
\focus{PC6}{Collect your group's result for the class.}
\focus{PA1}{Pool the class data in one table.}
\focus{PA2}{Plot a solubility curve.}
\focus{PA5}{Describe the trend.}
\focus{PA7}{Find and explain any anomalies.}
\focus{EV2}{Compare the class curve with a published one.}
\end{skillsfocus}

% ─── TECHNIQUE ────────────────────────────────────────────────────────────────
\exppart{Technique spotlight}{Blue}

\begin{technique}{Reaching saturation}{PC4}
Add the solid a little at a time and stir until each portion has dissolved. The solution is \textbf{saturated} when a portion will not dissolve, even after a full minute of stirring. Count only the portions that dissolved completely.
\begin{enumerate}[leftmargin=1.6em, itemsep=1pt, topsep=3pt]
  \item Weigh each portion on the balance to the nearest 0.1\,g.
  \item Keep the temperature steady in a water bath, and read the thermometer at eye level.
  \item Stir for the same time after every portion.
\end{enumerate}
\end{technique}

\begin{technique}{Plotting a line graph}{PA2}
\begin{enumerate}[leftmargin=1.6em, itemsep=1pt, topsep=3pt]
  \item Put the independent variable (temperature) on the horizontal axis and the dependent variable (mass dissolved) on the vertical axis. Give each a label and a unit.
  \item Choose an even scale that uses most of the grid: each square must stand for the same amount.
  \item Plot each point as a small neat cross ($\times$).
  \item Draw one smooth \textbf{line of best fit} through the middle of the points. Do not join the dots.
\end{enumerate}
\end{technique}

\begin{workedex}{Choosing a scale}
The largest mass is 10.5\,g and the grid is 16 squares high. 16 squares $\times$ 1\,g would fit, but 2 squares per gram (0.5\,g a square) would not (that needs 21 squares). Use 1\,g per square.
\end{workedex}

\begin{wedo}{Choosing a scale}
The temperatures go from 20\,\textdegree C to 60\,\textdegree C and the grid is 20 squares wide. Choose a scale for the horizontal axis and explain your choice.
\answerbox{3cm}
\end{wedo}

% ─── PLAN AND SAFETY / METHOD ────────────────────────────────────────────────
\exppart{Plan, safety and method}{Teal}

\textbf{Aim:} To find out how the temperature of water affects the mass of potassium nitrate that dissolves in it.

Each group tests one temperature, then the class pools its results.\quad Our temperature: \rule{2cm}{0.4pt}\,\textdegree C

\task Predict how the mass that dissolves will change as the water gets hotter. Explain using the particle model. \skr{QP1}
\answerbox{4cm}

\subsubsection*{Equipment}
Boiling tube \textbullet{} 10\,mL measuring cylinder \textbullet{} electronic balance (0.1\,g) \textbullet{} weighing boats \textbullet{} spatula \textbullet{} stirring rod \textbullet{} thermometer \textbullet{} 250\,mL beaker as a water bath, with kettle water mixed to your temperature \textbullet{} potassium nitrate

\begin{center}
\begin{tikzpicture}[scale=0.8]
  \eqBeaker
  \fill[black!8] (0.02,0.02) rectangle (1.78,1.9); \draw[Ink, line width=0.8pt] (0.02,1.9) -- (1.78,1.9);
  \begin{scope}[shift={(0.6,0.35)}]\eqTestTube\end{scope}
  \fill[black!25] (0.62,0.65) arc[start angle=180,end angle=360,radius=0.28] -- (1.18,1.6) -- (0.62,1.6) -- cycle;
  \begin{scope}[shift={(1.35,0.5)}, scale=0.95]\eqThermometer\end{scope}
  \eqlab{(1.2,3.3)}{(3.0,3.6)}{boiling tube}
  \eqlab{(1.6,2.9)}{(3.0,2.9)}{thermometer}
  \eqlab{(0.9,1.2)}{(3.0,1.8)}{10\,mL water $+$ potassium nitrate}
  \eqlab{(1.78,0.6)}{(3.0,0.6)}{water bath at your temperature}
\end{tikzpicture}
\expfigcap{The boiling tube stands in a water bath to hold its temperature.}
\end{center}

\begin{infobox}{Safety}
Wear safety glasses. Potassium nitrate is an oxidiser: keep it away from flames, and wash your hands after the lesson. Hot water scalds: your teacher fills the water bath. Put leftover solutions in the labelled waste container, not down the sink.
\end{infobox}

\subsubsection*{Method}
\begin{enumerate}[itemsep=2pt]
  \item Measure 10\,mL of water into the boiling tube and stand it in the water bath. Wait until the thermometer reads your temperature.
  \item Weigh 0.5\,g of potassium nitrate and add it to the tube. Stir until it has all dissolved.
  \item Keep adding 0.5\,g portions, one at a time, stirring after each.
  \item Stop when a portion will not dissolve after a minute of stirring. Record the total mass that \emph{did} dissolve.
  \item Write your result on the class table on the board, and copy the whole class's results into Table 5.1.
\end{enumerate}

% ─── RESULTS ─────────────────────────────────────────────────────────────────
\exppart{Results}{Amber}

\subsubsection*{Table 5.1: Mass of potassium nitrate that dissolves in 10\,mL of water (class data) \skr{PA1}}
\begin{ctab}
\begin{tabular}{|SW{3.0cm}|C{1.9cm}|C{1.9cm}|C{1.9cm}|C{1.9cm}|C{2.6cm}|}
  \hline\hd & \thead{Temperature (\textdegree C)} & \thead{Group 1 (g)} & \thead{Group 2 (g)} & \thead{Group 3 (g)} & \thead{Group 4 (g)} & \thead{Average (g)} \\ \hline
  & \scaf{20} & \scaf{3.0} & \scaf{3.5} & \scaf{3.0} & \scaf{---} & \scaf{3.2} \\ \hline
  & 20 & & & & & \\ \hline
  & 30 & & & & & \\ \hline
  & 40 & & & & & \\ \hline
  & 50 & & & & & \\ \hline
  & 60 & & & & & \\ \hline
\end{tabular}
\end{ctab}

\needspace{20\baselineskip}
\task Plot the class averages as a line graph and draw a smooth line of best fit. This is a \textbf{solubility curve}. \skr{PA2}
\graphgrid{10cm}

% ─── ANALYSIS ────────────────────────────────────────────────────────────────
\exppart{Analysis and evaluation}{Coral}

\task Describe the trend shown by your solubility curve. \skr{PA5}
\answerbox{4cm}

\needspace{10\baselineskip}
\task Look at the results of each group in Table 5.1. Is there an anomaly? Circle it and suggest what might have caused it. \skr{PA7}
\answerbox{4cm}

\needspace{10\baselineskip}
\task Use your curve to estimate the mass of potassium nitrate that would dissolve in 10\,mL of water at 45\,\textdegree C. Show on the graph how you found it.
\answerbox{2cm}

\needspace{14\baselineskip}
\task The published solubility of potassium nitrate is shown below. Compare it with the class results. Are the class values higher or lower? Suggest a reason. \skr{EV2}
\begin{ctab}
\begin{tabular}{|L{5.4cm}|C{1.6cm}|C{1.6cm}|C{1.6cm}|C{1.6cm}|C{1.6cm}|}
  \hline\thead{Temperature (\textdegree C)} & 20 & 30 & 40 & 50 & 60 \\ \hline
  \thead{Mass in 10\,mL of water (g)} & 3.2 & 4.6 & 6.4 & 8.5 & 11.0 \\ \hline
\end{tabular}
\end{ctab}
\answerbox{4cm}

\needspace{10\baselineskip}
\task A saturated solution at 60\,\textdegree C is cooled to 20\,\textdegree C. Predict what you would see, and explain using your curve.
\answerbox{4cm}

\begin{yr8box}
\begin{keepext}
\exttask The solubility of table salt hardly changes between 0\,\textdegree C and 100\,\textdegree C. Sketch on your graph what its solubility curve would look like, and explain why cooling a salt solution is not a good way to get salt crystals.
\answerbox{5cm}
\end{keepext}
\end{yr8box}

% ─── SKILL CHECK ─────────────────────────────────────────────────────────────
\exppart{Skill check}{NavyBlue}

\begin{skillcheck}
\scrow{PC4}{I measured mass to 0.1\,g and kept the temperature steady.}
\scrow{PC6}{I collected a result the class could use.}
\scrow{PA1}{I pooled the class data and calculated averages.}
\scrow{PA2}{I plotted a line graph with a line of best fit.}
\scrow{PA5}{I can describe the trend in a solubility curve.}
\scrow{PA7}{I can find an anomaly and suggest a cause.}
\scrow{EV2}{I can compare class results with published values.}
\end{skillcheck}

\task Which skill will you focus on in Experiment 6? Explain what you will do differently.
\answerbox{5cm}
% ══════════════════════════════════════════════════════════════════════════════
%  EXPERIMENT 6 — GROWING COPPER SULFATE CRYSTALS               (Lesson 14, Term 2 Week 5)
% ══════════════════════════════════════════════════════════════════════════════
\experimentdiv{Growing Copper Sulfate Crystals}{Handling a hazardous chemical \textbullet{} Safe disposal \textbullet{} Measuring growth over two weeks}{Lesson 14 \textbullet\ Term 2, Week 5 \textbullet\ Stretch lab}\label{sec:6}
\lessoncodes{VC2S8U05, U07 \textbullet{} VC2S8I02, I07, I08}

\vocabhead
\begin{ctab}
\begin{tabular}{|SW{4.2cm}|W{11.6cm}|}
  \hline\hd & \thead{Word} & \thead{Definition (copy from the slide)} \\ \hline
  & \scaf{\textbf{Seed crystal}\label{voc:seed-crystal}} & \scaf{A small crystal that the dissolved particles build onto as the solution evaporates.} \\ \hline
  & \textbf{Hazard}\label{voc:hazard} & \\ \hline
  & \textbf{Risk}\label{voc:risk} & \\ \hline
  & \textbf{Disposal}\label{voc:disposal} & \\ \hline
  & \textbf{Ethical}\label{voc:ethical} & \\ \hline
\end{tabular}
\end{ctab}

\begin{lipanel}
\inquiryli{Identify hazards and plan safe handling and disposal; record observations over time.}
\understandli{Explain how a crystal grows as the solvent evaporates from a saturated solution.}
\sctitle
\begin{itemize}[leftmargin=1.4em, itemsep=0pt, topsep=1pt]
  \item I can name the hazards of copper sulfate and how to control them.
  \item I can explain why copper sulfate must not go down the sink.
  \item I can measure and draw my crystal as it grows.
\end{itemize}
\end{lipanel}

\begin{skillsfocus}
\focus{PC3}{Name the hazards and control the risks.}
\focus{PC5}{Dispose of chemicals responsibly.}
\focus{PC6}{Measure the crystal every few days.}
\focus{CO2}{Draw the set-up and the crystal.}
\focus{EV5}{Suggest how to grow a bigger crystal.}
\end{skillsfocus}

% ─── TECHNIQUE ────────────────────────────────────────────────────────────────
\exppart{Technique spotlight}{Blue}

\begin{technique}{Hazard, risk and control}{PC3}
A \textbf{hazard} is something that could cause harm. The \textbf{risk} is how likely the harm is and how bad it would be. A \textbf{control} is what you do to make the risk smaller.
\end{technique}

\begin{workedex}{One row of a risk assessment}
\textbf{Hazard:} hot copper sulfate solution.\quad \textbf{Risk:} it could scald skin if spilled.\quad \textbf{Control:} the teacher pours it; students keep beakers in the middle of the bench.
\end{workedex}

\begin{youdobox}[Risk assessment]
\task Complete the risk assessment for this experiment. \skr{PC3}
\begin{ctab}
\begin{tabular}{|SW{3.8cm}|W{4.6cm}|W{5.8cm}|}
  \hline\hd & \thead{Hazard} & \thead{Risk} & \thead{Control} \\ \hline
  & \scaf{Copper sulfate is harmful if swallowed} & \scaf{Could make someone ill} & \scaf{Never taste; wash hands after handling} \\ \hline
  & Copper sulfate irritates eyes and skin & & \\ \hline
  & Glass beaker & & \\ \hline
  & Hot solution & & \\ \hline
\end{tabular}
\end{ctab}
\end{youdobox}

\begin{technique}{Responsible disposal}{PC5}
Copper sulfate is \textbf{toxic to fish and other water life}. Water from sinks flows to the sewage plant and on to rivers and the bay, and the treatment plant does not remove all of it. Scientists have a responsibility to the environment as well as to themselves: every drop of copper sulfate solution goes into the labelled \textbf{heavy metal waste} container, never down the sink. Crystals you keep stay in a sealed, labelled bag.
\end{technique}

\begin{wedo}{Why it matters}
A student says, ``It's only a few drops, it won't matter if I tip it down the sink.'' Explain why the class still uses the waste container.
\answerbox{3.5cm}
\end{wedo}

% ─── PLAN AND SAFETY / METHOD ────────────────────────────────────────────────
\exppart{Plan, safety and method}{Teal}

\textbf{Aim:} To grow a large crystal of copper sulfate from a seed crystal.

\subsubsection*{Equipment}
100\,mL beaker \textbullet{} saturated copper sulfate solution (made hot by your teacher) \textbullet{} a small seed crystal \textbullet{} cotton thread \textbullet{} pencil or paddle-pop stick \textbullet{} paper towel and rubber band \textbullet{} ruler \textbullet{} hand lens \textbullet{} gloves and safety glasses

\begin{center}
\begin{tikzpicture}[scale=0.8]
  \eqBeaker
  \fill[black!15] (0.02,0.02) rectangle (1.78,1.8); \draw[Ink, line width=0.8pt] (0.02,1.8) -- (1.78,1.8);
  \draw[eq] (-0.5,2.75) -- (2.3,2.75);
  \draw[Ink, line width=0.8pt] (0.9,2.75) -- (0.9,1.0);
  \draw[eq, fill=black!55] (0.9,1.05) -- (1.1,0.85) -- (1.0,0.6) -- (0.8,0.6) -- (0.7,0.85) -- cycle;
  \eqlab{(2.2,2.75)}{(3.2,3.1)}{pencil across the top}
  \eqlab{(0.9,1.4)}{(3.2,2.2)}{cotton thread}
  \eqlab{(1.1,0.85)}{(3.2,1.2)}{seed crystal, off the bottom}
  \eqlab{(1.78,0.4)}{(3.2,0.4)}{saturated copper sulfate solution}
\end{tikzpicture}
\expfigcap{The seed crystal hangs in the solution without touching the glass.}
\end{center}

\begin{infobox}{Safety}
Wear safety glasses and gloves. Copper sulfate is harmful if swallowed and irritates eyes and skin: never taste it, and wash your hands after the lesson even if you wore gloves. Your teacher pours the hot solution. All leftover solution goes in the heavy metal waste container.
\end{infobox}

\subsubsection*{Method}
\begin{enumerate}[itemsep=2pt]
  \item Tie the seed crystal to one end of the thread and the other end to the pencil.
  \item Your teacher half fills your beaker with warm saturated solution.
  \item Lay the pencil across the beaker so that the crystal hangs in the middle, clear of the bottom and the sides.
  \item Cover the beaker with paper towel held on by a rubber band, to keep dust out but let water evaporate. Label it.
  \item Leave it somewhere still. Measure and draw the crystal at each lesson for two weeks.
  \item At the end, put the leftover solution in the waste container and dry the crystal on paper towel.
\end{enumerate}

% ─── RESULTS ─────────────────────────────────────────────────────────────────
\exppart{Results}{Amber}

\subsubsection*{Table 6.1: Growth of the copper sulfate crystal \skr{PC6}}
\begin{ctab}
\begin{tabular}{|SW{2.6cm}|C{2.6cm}|W{9.8cm}|}
  \hline\hd & \thead{Day} & \thead{Longest side (mm)} & \thead{Observations (colour, shape, crystals on the glass)} \\ \hline
  & \scaf{Example} & \scaf{3} & \scaf{Blue, four-sided; a crust forming at the water line} \\ \hline
  & Day 0 & & \\ \hline
  & Day 2 & & \\ \hline
  & Day 5 & & \\ \hline
  & Day 7 & & \\ \hline
  & Day 10 & & \\ \hline
  & Day 14 & & \\ \hline
\end{tabular}
\end{ctab}

\needspace{14\baselineskip}
\task Draw your crystal through the hand lens on the first and last days. Label its faces and edges, and write its length under each drawing. \skr{CO2}

\noindent
\begin{minipage}[t]{0.48\linewidth}\drawbox{6cm}{}\nameline{Day 0}\end{minipage}\hfill
\begin{minipage}[t]{0.48\linewidth}\drawbox{6cm}{}\nameline{Day 14}\end{minipage}

% ─── ANALYSIS ────────────────────────────────────────────────────────────────
\exppart{Analysis and evaluation}{Coral}

\task Explain how the crystal grew. Use the words \emph{saturated}, \emph{evaporate} and \emph{particles}.
\answerbox{5cm}

\needspace{10\baselineskip}
\task Small crystals often formed on the sides of the beaker too. Suggest why, and why they take copper sulfate away from your seed crystal.
\answerbox{4cm}

\needspace{10\baselineskip}
\task Why must the solution be saturated before the seed crystal goes in? What would happen to the seed crystal if it were not?
\answerbox{4cm}

\needspace{10\baselineskip}
\task Explain why the leftover solution went into the waste container and not down the sink. \skr{PC5}
\answerbox{4cm}

\needspace{10\baselineskip}
\task Suggest one change that could grow a larger crystal, and write it as a question you could investigate. \skr{EV5}
\answerbox{4cm}

\begin{yr8box}
\begin{keepext}
\exttask Copper sulfate crystals are blue, but if you heat them they turn white and give off steam. Suggest what the crystals contain besides copper sulfate, and how you could test your idea.
\answerbox{5cm}
\end{keepext}
\end{yr8box}

% ─── SKILL CHECK ─────────────────────────────────────────────────────────────
\exppart{Skill check}{NavyBlue}

\begin{skillcheck}
\scrow{PC3}{I can name the hazards of copper sulfate and how to control them.}
\scrow{PC5}{I disposed of the solution responsibly and can explain why.}
\scrow{PC6}{I measured and recorded my crystal at each lesson.}
\scrow{CO2}{I drew my crystal clearly and labelled it.}
\scrow{EV5}{I can suggest a question for a further investigation.}
\end{skillcheck}

\task Which skill will you focus on in Experiment 7? Explain what you will do differently.
\answerbox{5cm}
% ══════════════════════════════════════════════════════════════════════════════
%  EXPERIMENT 7 — FILTRATION                                     (Lesson 15, Term 2 Week 6)
% ══════════════════════════════════════════════════════════════════════════════
\experimentdiv{Filtration}{Folding and using a filter paper \textbullet{} Comparing filters fairly \textbullet{} Residue and filtrate}{Lesson 15 \textbullet\ Term 2, Week 6 \textbullet\ Lab}\label{sec:7}
\lessoncodes{VC2S8U06, U07 \textbullet{} VC2S8I02, I03, I06, I08}

\vocabhead
\begin{ctab}
\begin{tabular}{|SW{4.2cm}|W{11.6cm}|}
  \hline\hd & \thead{Word} & \thead{Definition (copy from the slide)} \\ \hline
  & \scaf{\textbf{Filtration}\label{voc:filtration}} & \scaf{Separating an insoluble solid from a liquid by passing it through a filter.} \\ \hline
  & \textbf{Filtrate}\label{voc:filtrate} & \\ \hline
  & \textbf{Residue} & \\ \hline
  & \textbf{Decanting}\label{voc:decanting} & \\ \hline
  & \textbf{Pore}\label{voc:pore} & \\ \hline
\end{tabular}
\end{ctab}

\begin{lipanel}
\inquiryli{Set up filtration correctly and compare three filters in a fair test.}
\understandli{Explain how a filter separates an insoluble solid from a liquid.}
\sctitle
\begin{itemize}[leftmargin=1.4em, itemsep=0pt, topsep=1pt]
  \item I can fold a filter paper and set up a funnel correctly.
  \item I can name the residue and the filtrate.
  \item I can compare filters and explain which works best.
\end{itemize}
\end{lipanel}

\begin{skillsfocus}
\focus{PC2}{Keep the test of each filter fair.}
\focus{PC7}{Set up the filtration equipment correctly.}
\focus{PC6}{Time each filter and rate the filtrate.}
\focus{CO2}{Draw the set-up as a scientific diagram.}
\focus{EV3}{Find sources of error.}
\end{skillsfocus}

% ─── TECHNIQUE ────────────────────────────────────────────────────────────────
\exppart{Technique spotlight}{Blue}

\begin{technique}{Folding a filter paper}{PC7}
\begin{center}
\begin{tikzpicture}[scale=0.62]
  \draw[eq] (0,0) circle (1.2);
  \draw[Ink, dashed, line width=0.7pt] (-1.2,0) -- (1.2,0);
  \node[eqtext, align=center] at (0,-1.9) {1. Fold\\in half};
  \begin{scope}[shift={(4.2,-0.3)}]
    \draw[eq] (-1.2,0) arc[start angle=180,end angle=0,radius=1.2] -- cycle;
    \draw[Ink, dashed, line width=0.7pt] (0,0) -- (0,1.2);
    \node[eqtext, align=center] at (0,-1.6) {2. Fold in\\half again};
  \end{scope}
  \begin{scope}[shift={(8.4,-0.3)}]
    \draw[eq] (0,0) -- (1.2,0) arc[start angle=0,end angle=90,radius=1.2] -- cycle;
    \node[eqtext, align=center] at (0.5,-1.6) {3. Open one layer\\to make a cone};
  \end{scope}
  \begin{scope}[shift={(13.0,-1.1)}]
    \draw[eq] (0,2.4) -- (0.9,0.9) -- (0.9,0) (2.2,2.4) -- (1.3,0.9) -- (1.3,0);
    \draw[Ink, line width=1.1pt, fill=black!10] (0.25,2.25) -- (1.1,1.05) -- (1.95,2.25) -- cycle;
    \node[eqtext, align=center] at (1.1,-0.8) {4. Fit it in the funnel;\\wet it to stick};
  \end{scope}
\end{tikzpicture}
\end{center}
\begin{itemize}[leftmargin=1.4em, itemsep=1pt, topsep=2pt]
  \item Keep the top of the paper \textbf{below} the rim of the funnel, and the liquid \textbf{below} the top of the paper, or the mixture runs round the outside.
  \item Pour down a stirring rod so that the liquid does not splash or tear the paper.
  \item Never poke or stir inside the filter paper.
\end{itemize}
\end{technique}

\begin{workedex}{Residue and filtrate}
Muddy water is filtered. The \textbf{residue} is the soil left on the filter paper. The \textbf{filtrate} is the clearer water that passes through.
\end{workedex}

\begin{wedo}{Residue and filtrate}
Sand and salt water are filtered. Name the residue and the filtrate, and explain why the salt does not stay on the filter paper.
\answerbox{4cm}
\end{wedo}

% ─── PLAN AND SAFETY / METHOD ────────────────────────────────────────────────
\exppart{Plan, safety and method}{Teal}

\textbf{Aim:} To compare how well filter paper, a coffee filter and a cotton cloth filter muddy water.

\task Name the variables. \skr{PC2}
\begin{ctab}
\begin{tabular}{|SW{4.4cm}|W{10.8cm}|}
  \hline\hd & \thead{Variable} & \thead{In this experiment} \\ \hline
  & \scaf{Independent} & \scaf{The type of filter} \\ \hline
  & Dependent (two) & \\ \hline
  & Controlled & \\ \hline
  & Controlled & \\ \hline
\end{tabular}
\end{ctab}

\needspace{8\baselineskip}
\task Predict which filter will give the clearest filtrate, and which will be fastest. Give a reason for each. \skr{QP1}
\answerbox{2.4cm}

\subsubsection*{Equipment}
Filter funnel \textbullet{} retort stand and ring clamp (or a conical flask to sit the funnel in) \textbullet{} 3 beakers \textbullet{} 100\,mL measuring cylinder \textbullet{} stirring rod \textbullet{} filter paper, coffee filter and a square of cotton cloth \textbullet{} 300\,mL of muddy water (stir well before each pour) \textbullet{} stopwatch \textbullet{} white card

\begin{infobox}{Safety}
Wear safety glasses. Glass funnels and beakers break easily: keep them in the middle of the bench. Pour the muddy residue into the bucket, not the sink, so the drain does not block.
\end{infobox}

\subsubsection*{Method}
\begin{enumerate}[itemsep=2pt]
  \item Fold the filter paper and fit it in the funnel. Stand the funnel in a clean beaker.
  \item Stir the muddy water, then measure 50\,mL of it.
  \item Pour it into the filter down the stirring rod and start the stopwatch.
  \item Stop the stopwatch when drips are more than 5 seconds apart. Record the time and the volume of filtrate.
  \item Hold the beaker against the white card and rate the filtrate's clearness from 1 (as muddy as at the start) to 5 (completely clear).
  \item Repeat with the coffee filter and then the cloth, using a clean beaker each time.
\end{enumerate}

% ─── RESULTS ─────────────────────────────────────────────────────────────────
\exppart{Results}{Amber}

\subsubsection*{Table 7.1: Comparing three filters \skr{PC6}}
\begin{ctab}
\begin{tabular}{|SW{3.4cm}|C{2.4cm}|C{2.8cm}|C{2.6cm}|W{3.2cm}|}
  \hline\hd & \thead{Filter} & \thead{Time (s)} & \thead{Volume of filtrate (mL)} & \thead{Clearness (1--5)} & \thead{Residue} \\ \hline
  & \scaf{Paper towel} & \scaf{95} & \scaf{38} & \scaf{3} & \scaf{Thin brown layer} \\ \hline
  & Filter paper & & & & \\ \hline
  & Coffee filter & & & & \\ \hline
  & Cotton cloth & & & & \\ \hline
\end{tabular}
\end{ctab}

\needspace{20\baselineskip}
\task Draw a scientific diagram of your filtration set-up. Label the funnel, filter paper, residue, filtrate and beaker. \skr{CO2}
\drawbox{8cm}{}

% ─── ANALYSIS ────────────────────────────────────────────────────────────────
\exppart{Analysis and evaluation}{Coral}

\task Which filter gave the clearest filtrate, and which was fastest? Use your results.
\answerbox{4cm}

\needspace{12\baselineskip}
\task The clearest filter was not the fastest. Explain why, using the idea of the size of the holes (pores) in the filter.
\answerbox{4.5cm}

\needspace{10\baselineskip}
\task Even the clearest filtrate may not be safe to drink. Explain why. What could be done to make it safe?
\answerbox{4cm}

\needspace{12\baselineskip}
\task Give two sources of error and how each could be reduced. \skr{EV3}
\begin{ctab}
\begin{tabular}{|SW{7.2cm}|W{7.8cm}|}
  \hline\hd & \thead{Source of error} & \thead{How to reduce it} \\ \hline
  & \scaf{The mud settled, so later pours were less muddy} & \scaf{Stir the muddy water for the same time before every pour} \\ \hline
  & & \\ \hline
  & & \\ \hline
\end{tabular}
\end{ctab}

\needspace{10\baselineskip}
\task Decanting is quicker than filtering. Explain when you would decant first and then filter.
\answerbox{4cm}

\begin{yr8box}
\begin{keepext}
\exttask Melbourne's drinking water passes through layers of sand and gravel at a treatment plant. Explain how a sand filter works, and why it has gravel at the bottom and fine sand at the top.
\answerbox{5cm}
\end{keepext}
\end{yr8box}

% ─── SKILL CHECK ─────────────────────────────────────────────────────────────
\exppart{Skill check}{NavyBlue}

\begin{skillcheck}
\scrow{PC2}{I kept the test of each filter fair.}
\scrow{PC7}{I folded the filter paper and set up the funnel correctly.}
\scrow{PC6}{I timed each filter and rated each filtrate.}
\scrow{CO2}{I drew and labelled the filtration set-up.}
\scrow{EV3}{I can name sources of error and how to reduce them.}
\end{skillcheck}

\task Which skill will you focus on in Experiment 8? Explain what you will do differently.
\answerbox{5cm}
% ══════════════════════════════════════════════════════════════════════════════
%  EXPERIMENT 8 — DISTILLATION                                   (Lesson 17, Term 2 Week 6)
% ══════════════════════════════════════════════════════════════════════════════
\experimentdiv{Distillation}{Setting up a distillation apparatus \textbullet{} Heating a liquid safely \textbullet{} Testing the distillate}{Lesson 17 \textbullet\ Term 2, Week 6 \textbullet\ Lab}\label{sec:8}
\lessoncodes{VC2S8U05, U07 \textbullet{} VC2S8I02, I03, I06, I08}

\vocabhead
\begin{ctab}
\begin{tabular}{|SW{4.2cm}|W{11.6cm}|}
  \hline\hd & \thead{Word} & \thead{Definition (copy from the slide)} \\ \hline
  & \scaf{\textbf{Distillation}\label{voc:distillation}} & \scaf{Separating a liquid from a solution by boiling it and condensing the vapour.} \\ \hline
  & \textbf{Distillate}\label{voc:distillate} & \\ \hline
  & \textbf{Condense}\label{voc:condense} & \\ \hline
  & \textbf{Boiling point}\label{voc:boiling-point} & \\ \hline
  & \textbf{Boiling chip}\label{voc:boiling-chip} & \\ \hline
\end{tabular}
\end{ctab}

\begin{lipanel}
\inquiryli{Select and set up equipment for distillation, and use it safely.}
\understandli{Explain how distillation separates pure water from a solution, using changes of state.}
\sctitle
\begin{itemize}[leftmargin=1.4em, itemsep=0pt, topsep=1pt]
  \item I can name each piece of equipment and say what it does.
  \item I can draw and label the distillation set-up.
  \item I can explain where the water went and what was left behind.
\end{itemize}
\end{lipanel}

\begin{skillsfocus}
\focus{PC7}{Choose and set up the distillation equipment.}
\focus{PC3}{Heat a flask safely.}
\focus{PC6}{Test the distillate and the residue.}
\focus{CO2}{Draw the set-up as a scientific diagram.}
\focus{EV2}{Compare the results with the water cycle.}
\end{skillsfocus}

% ─── TECHNIQUE ────────────────────────────────────────────────────────────────
\exppart{Technique spotlight}{Blue}

\begin{technique}{The distillation set-up}{PC7}
\begin{center}
\begin{tikzpicture}[scale=0.78]
  % stand and clamp holding the flask
  \begin{scope}[shift={(-1.6,0)}]\draw[eq] (0,0) rectangle (2.0,0.2); \draw[eq] (0.3,0.2) -- (0.3,5.6) (0.45,0.2) -- (0.45,5.6);\end{scope}
  \draw[eq] (-1.15,4.6) -- (0.6,4.6) (-1.15,4.75) -- (0.6,4.75);
  % tripod, gauze, Bunsen
  \begin{scope}[shift={(-0.35,0.2)}]\eqBunsen\end{scope}
  \begin{scope}[shift={(-0.6,2.2)}]\draw[eq] (-0.15,0.0) rectangle (2.55,0.1);\draw[eq] (0.25,0) -- (0,-1.95) (2.15,0) -- (2.4,-1.95);\end{scope}
  % conical flask
  \begin{scope}[shift={(-0.15,2.3)}]\eqFlask\end{scope}
  \fill[black!15] (-0.12,2.33) -- (0.22,3.3) -- (1.53,3.3) -- (1.72,2.33) -- cycle;
  \draw[Ink, line width=0.8pt] (0.22,3.3) -- (1.53,3.3);
  \foreach \p in {(0.5,2.5),(0.9,2.6),(1.2,2.45)}{\draw[Ink, fill=white] \p circle (0.07);}
  \fill[Ink] (0.45,5.1) rectangle (1.15,5.45);
  % delivery tube
  \draw[Ink, line width=1.6pt] (0.8,5.45) -- (0.8,5.8) -- (4.4,4.3) -- (4.4,2.0);
  % ice bath and collecting tube
  \begin{scope}[shift={(3.5,0.2)}]\eqBeaker\end{scope}
  \fill[black!6] (3.52,0.22) rectangle (5.28,2.0); \draw[Ink, line width=0.8pt] (3.52,2.0) -- (5.28,2.0);
  \foreach \p in {(3.8,1.6),(5.0,1.2),(3.9,0.7),(5.0,0.5)}{\draw[Ink] \p ++(-0.15,-0.15) rectangle ++(0.3,0.3);}
  \begin{scope}[shift={(4.1,0.4)}]\eqTestTube\end{scope}
  \fill[black!3] (4.12,0.7) arc[start angle=180,end angle=360,radius=0.28] -- (4.68,1.0) -- (4.12,1.0) -- cycle;
  \eqlabl{(0.2,3.0)}{(-2.2,3.6)}{salty coloured water}
  \eqlabl{(0.5,2.5)}{(-2.2,2.6)}{boiling chips}
  \eqlabl{(0.45,5.3)}{(-2.2,5.9)}{rubber bung}
  \eqlab{(2.6,5.05)}{(3.6,5.9)}{delivery tube}
  \eqlab{(5.28,1.5)}{(6.0,1.9)}{ice and water}
  \eqlab{(4.68,0.85)}{(6.0,0.9)}{distillate}
  \eqlabl{(-0.1,0.8)}{(-2.2,1.4)}{Bunsen burner}
\end{tikzpicture}
\end{center}
\begin{itemize}[leftmargin=1.4em, itemsep=1pt, topsep=2pt]
  \item The \textbf{delivery tube} slopes \emph{down} to the collecting tube, so the distillate runs away from the flask.
  \item The collecting tube sits in \textbf{ice water}, which cools the vapour so that it condenses.
  \item \textbf{Boiling chips} make the liquid boil smoothly instead of in sudden violent bursts.
  \item The end of the delivery tube stays \textbf{above} the liquid in the collecting tube, or cold water can be sucked back into the hot flask.
\end{itemize}
\end{technique}

\begin{youdobox}[Choosing equipment]
\task Complete the table. \skr{PC7}
\begin{ctab}
\begin{tabular}{|SW{4.0cm}|W{11.0cm}|}
  \hline\hd & \thead{Equipment} & \thead{Its job in distillation} \\ \hline
  & \scaf{Bunsen burner} & \scaf{Heats the solution until the water boils} \\ \hline
  & Conical flask & \\ \hline
  & Rubber bung & \\ \hline
  & Delivery tube & \\ \hline
  & Beaker of ice water & \\ \hline
  & Boiling chips & \\ \hline
\end{tabular}
\end{ctab}
\end{youdobox}

% ─── PLAN AND SAFETY / METHOD ────────────────────────────────────────────────
\exppart{Plan, safety and method}{Teal}

\textbf{Aim:} To separate pure water from salty, coloured water by distillation.

\task Predict the colour of the distillate and what will be left in the flask. Explain your prediction. \skr{QP1}
\answerbox{4cm}

\subsubsection*{Equipment}
250\,mL conical flask \textbullet{} rubber bung with a bent delivery tube \textbullet{} retort stand and clamp \textbullet{} tripod, gauze mat, Bunsen burner and heatproof mat \textbullet{} test tube \textbullet{} 250\,mL beaker of ice and water \textbullet{} 3 boiling chips \textbullet{} 50\,mL of salt water with blue food colouring \textbullet{} 2 watch glasses \textbullet{} dropper

\begin{infobox}{Safety}
Wear safety glasses. Your teacher checks the set-up before you light the Bunsen. Never heat a sealed container: the delivery tube must be open at the end. Turn the gas off \emph{before} the flask boils dry, and lift the delivery tube out of the collecting tube when you stop heating. The flask, tube and steam stay hot for several minutes.
\end{infobox}

\subsubsection*{Method}
\begin{enumerate}[itemsep=2pt]
  \item Put 50\,mL of salty coloured water and three boiling chips into the conical flask.
  \item Clamp the flask above the tripod and gauze, and fit the bung and delivery tube.
  \item Stand the test tube in the ice water with the end of the delivery tube inside it, above the bottom.
  \item Ask your teacher to check the set-up, then heat gently with a blue flame.
  \item When the test tube is about a quarter full, lift the delivery tube out, then turn off the gas.
  \item Put one drop of the distillate on a watch glass and one drop of the original solution on another. Leave them on the windowsill until next lesson, then compare.
\end{enumerate}

% ─── RESULTS ─────────────────────────────────────────────────────────────────
\exppart{Results}{Amber}

\subsubsection*{Table 8.1: Observations during distillation \skr{PC6}}
\begin{ctab}
\begin{tabular}{|SW{5.0cm}|W{10.0cm}|}
  \hline\hd & \thead{What was observed} & \thead{Observation} \\ \hline
  & \scaf{Solution in the flask at the start} & \scaf{Clear, blue liquid} \\ \hline
  & Inside the flask while it boiled & \\ \hline
  & Inside the delivery tube & \\ \hline
  & Distillate in the test tube & \\ \hline
  & Liquid left in the flask & \\ \hline
  & Watch glass: drop of distillate (next lesson) & \\ \hline
  & Watch glass: drop of original solution & \\ \hline
\end{tabular}
\end{ctab}

\needspace{20\baselineskip}
\task Draw a scientific diagram of your set-up without looking at the technique box. Label every piece of equipment, the solution and the distillate. \skr{CO2}
\drawbox{9cm}{}

% ─── ANALYSIS ────────────────────────────────────────────────────────────────
\exppart{Analysis and evaluation}{Coral}

\task Describe what happened to the water, naming the two changes of state.
\answerbox{4cm}

\needspace{10\baselineskip}
\task How do the watch glasses show that the distillate was pure water? \skr{PC6}
\answerbox{4cm}

\needspace{10\baselineskip}
\task Why did the salt and colouring stay in the flask? Use the idea of boiling point.
\answerbox{4cm}

\needspace{10\baselineskip}
\task Compare distillation with the water cycle. Which part of the water cycle is like the Bunsen burner, and which is like the ice water? \skr{EV2}
\answerbox{4.5cm}

\needspace{10\baselineskip}
\task Never taste the distillate, even though it should be pure water. Give two reasons.
\answerbox{3.5cm}

\begin{yr8box}
\begin{keepext}
\exttask Wonthaggi's desalination plant uses reverse osmosis rather than distillation to make drinking water from seawater. Suggest why distillation would cost much more to run for a whole city.
\answerbox{4.5cm}
\end{keepext}
\end{yr8box}

% ─── SKILL CHECK ─────────────────────────────────────────────────────────────
\exppart{Skill check}{NavyBlue}

\begin{skillcheck}
\scrow{PC7}{I can name each piece of distillation equipment and its job.}
\scrow{PC3}{I heated the flask safely and stopped before it boiled dry.}
\scrow{PC6}{I tested the distillate and compared it with the solution.}
\scrow{CO2}{I can draw and label the set-up from memory.}
\scrow{EV2}{I can compare distillation with the water cycle.}
\end{skillcheck}

\task Which skill will you focus on in Experiment 9? Explain what you will do differently.
\answerbox{5cm}
% ══════════════════════════════════════════════════════════════════════════════
%  EXPERIMENT 9 — THE GREAT GREEN MIX UP                         (Lesson 18, Term 2 Week 7)
% ══════════════════════════════════════════════════════════════════════════════
% NOTE FOR KODIE: written without the faculty GGMU sheet, as a general
% "plan your own separation" practical. The mixture below (split peas, sand,
% iron filings, salt, green food colouring) is an assumption: swap it for the
% faculty mixture and the tasks still work.
\experimentdiv{The Great Green Mix Up}{Planning a sequence of separations \textbullet{} Writing a method in the third person \textbullet{} Judging how well it worked}{Lesson 18 \textbullet\ Term 2, Week 7 \textbullet\ Lab}\label{sec:9}
\lessoncodes{VC2S8U06, U07 \textbullet{} VC2S8I01--I03, I06, I08}

\vocabhead
\begin{ctab}
\begin{tabular}{|SW{4.2cm}|W{11.6cm}|}
  \hline\hd & \thead{Word} & \thead{Definition (copy from the slide)} \\ \hline
  & \scaf{\textbf{Mixture}\label{voc:mixture}} & \scaf{Two or more substances mixed together but not chemically joined.} \\ \hline
  & \textbf{Property}\label{voc:property} & \\ \hline
  & \textbf{Sieving}\label{voc:sieving} & \\ \hline
  & \textbf{Magnetic separation}\label{voc:magnetic-separation} & \\ \hline
  & \textbf{Flow chart}\label{voc:flow-chart} & \\ \hline
\end{tabular}
\end{ctab}

\begin{lipanel}
\inquiryli{Plan, carry out and evaluate a sequence of separations; write the method in the third person.}
\understandli{Choose a separation technique for each substance based on its properties.}
\sctitle
\begin{itemize}[leftmargin=1.4em, itemsep=0pt, topsep=1pt]
  \item I can match each substance to a property that lets me separate it.
  \item I can put the separations in an order that works, and show it as a flow chart.
  \item I can say how pure each separated substance was.
\end{itemize}
\end{lipanel}

\begin{skillsfocus}
\focus{QP2}{Write the aim of the challenge.}
\focus{PC1}{Plan the order of the separations.}
\focus{PC3}{Plan the safety of every step.}
\focus{PC7}{Choose equipment for each separation.}
\focus{CO1}{Write the method in the third person.}
\focus{EV3}{Judge how well each separation worked.}
\end{skillsfocus}

% ─── TECHNIQUE ────────────────────────────────────────────────────────────────
\exppart{Technique spotlight}{Blue}

\begin{technique}{Choosing a separation by property}{PC7}
Every separation uses a \textbf{property} that one substance has and the others do not.
\begin{ctab}
\renewcommand{\arraystretch}{1.3}
\begin{tabular}{|L{4.4cm}|L{4.6cm}|L{5.4cm}|}
  \hline\thead{Property} & \thead{Technique} & \thead{Example} \\ \hline
  Size of the pieces & Sieving & Gravel from sand \\ \hline
  Magnetic & Magnetic separation & Steel cans from aluminium cans \\ \hline
  Insoluble in water & Filtration, decanting & Sand from salt water \\ \hline
  Soluble in water & Dissolving, then evaporation & Salt from sand \\ \hline
  Different boiling points & Distillation & Pure water from salt water \\ \hline
  Different solubility in a solvent & Chromatography & The dyes in green ink \\ \hline
\end{tabular}
\end{ctab}
\end{technique}

\begin{workedex}{Planning the order}
\textbf{Mixture:} sand, salt and paperclips.
\begin{enumerate}[leftmargin=1.6em, itemsep=0pt, topsep=2pt]
  \item Magnetic separation removes the paperclips (steel is magnetic).
  \item Water is added and stirred: the salt dissolves, the sand does not.
  \item Filtration: the sand is the residue, the salt water is the filtrate.
  \item Evaporation of the filtrate leaves the salt.
\end{enumerate}
The magnet goes first: once the water is added, wet sand would stick to the paperclips.
\end{workedex}

\begin{wedo}{Planning the order}
A mixture contains rice, flour and salt. Plan the order of separations and give the property each step uses.
\answerbox{5cm}
\end{wedo}

% ─── PLAN ────────────────────────────────────────────────────────────────────
\exppart{Plan the separation}{Teal}

Your teacher gives your group a beaker of the \textbf{Great Green Mix Up}: green split peas, sand, iron filings and salt, all dyed with green food colouring. Your challenge is to get back as much of each as you can, as pure as you can.

\task Write the aim of the challenge. \skr{QP2}
\answerbox{2.5cm}

\needspace{12\baselineskip}
\task Complete the table. \skr{PC7}
\begin{ctab}
\begin{tabular}{|SW{3.4cm}|W{4.8cm}|W{3.6cm}|W{3.2cm}|}
  \hline\hd & \thead{Substance} & \thead{Property I will use} & \thead{Technique} & \thead{Equipment} \\ \hline
  & \scaf{Split peas} & \scaf{Much larger pieces} & \scaf{Sieving} & \scaf{Kitchen sieve} \\ \hline
  & Iron filings & & & \\ \hline
  & Sand & & & \\ \hline
  & Salt & & & \\ \hline
  & Green colouring & & & \\ \hline
\end{tabular}
\end{ctab}

\needspace{20\baselineskip}
\task Draw a flow chart of your separations, from the whole mixture at the top to the separate substances at the bottom. Write the technique on each arrow. \skr{PC1}
\drawbox{9cm}{}

\needspace{12\baselineskip}
\task Complete the risk assessment. \skr{PC3}
\begin{ctab}
\begin{tabular}{|SW{4.6cm}|W{5.0cm}|W{5.4cm}|}
  \hline\hd & \thead{Hazard} & \thead{Risk} & \thead{Control} \\ \hline
  & \scaf{Iron filings} & \scaf{Could get in eyes} & \scaf{Safety glasses; keep the magnet in a plastic bag} \\ \hline
  & & & \\ \hline
  & & & \\ \hline
\end{tabular}
\end{ctab}

{\small Teacher's initials before you start: \rule{2.5cm}{0.4pt}}

% ─── RESULTS ─────────────────────────────────────────────────────────────────
\exppart{Results}{Amber}

\task Write the method your group actually followed, as numbered steps in the third person. \skr{CO1}
\begin{ctab}
\begin{tabular}{|SW{1.4cm}|W{13.6cm}|}
  \hline\hd & \thead{Step} & \thead{What was done} \\ \hline
  & \scaf{1} & \scaf{The mixture was poured through a sieve and the peas were collected on paper towel.} \\ \hline
  & 2 & \\ \hline
  & 3 & \\ \hline
  & 4 & \\ \hline
  & 5 & \\ \hline
  & 6 & \\ \hline
  & 7 & \\ \hline
\end{tabular}
\end{ctab}

\needspace{12\baselineskip}
\subsubsection*{Table 9.1: How pure was each substance?}
\begin{ctab}
\begin{tabular}{|SW{3.0cm}|C{3.2cm}|W{8.6cm}|}
  \hline\hd & \thead{Substance} & \thead{Purity (1 to 5)} & \thead{What was still mixed in, or lost} \\ \hline
  & \scaf{Split peas} & \scaf{4} & \scaf{Still stained green; a little sand stuck to them} \\ \hline
  & Split peas & & \\ \hline
  & Iron filings & & \\ \hline
  & Sand & & \\ \hline
  & Salt & & \\ \hline
\end{tabular}
\end{ctab}

% ─── ANALYSIS ────────────────────────────────────────────────────────────────
\exppart{Analysis and evaluation}{Coral}

\task Which separation worked best? Explain why, using the property it relied on.
\answerbox{4cm}

\needspace{10\baselineskip}
\task Which substance was least pure? Name what was still mixed in, and how it got there. \skr{EV3}
\answerbox{4cm}

\needspace{10\baselineskip}
\task Did the order you planned work? If you changed it during the lesson, explain why.
\answerbox{4cm}

\needspace{10\baselineskip}
\task The green colouring dissolved in the water, so it ended up with the salt. Suggest how the salt could have been made white again, or how the colouring could be separated from the water.
\answerbox{4cm}

\needspace{10\baselineskip}
\task Suggest two improvements to your method and explain how each would make a substance purer. \skr{EV3}
\answerbox{5cm}

\begin{yr8box}
\begin{keepext}
\exttask Recycling plants separate mixed household waste. Choose three materials that a plant must separate, and for each give the property and technique it could use.
\answerbox{5cm}
\end{keepext}
\end{yr8box}

% ─── SKILL CHECK ─────────────────────────────────────────────────────────────
\exppart{Skill check}{NavyBlue}

\begin{skillcheck}
\scrow{QP2}{I can write the aim of a challenge.}
\scrow{PC1}{I planned the separations in an order that works.}
\scrow{PC3}{I planned how to make each step safe.}
\scrow{PC7}{I chose equipment that suits each separation.}
\scrow{CO1}{I wrote the method in the third person.}
\scrow{EV3}{I can judge how pure each substance was and suggest improvements.}
\end{skillcheck}

\task Look back over all nine experiments. Which skill have you improved most this unit? Give an example.
\answerbox{5cm}

\clearpage
\hypertarget{glossary}{}\pdfbookmark[1]{Glossary}{glossarybm}
\markboth{Glossary}{}{\headingfont\bfseries\LARGE Glossary}\label{sec:glossary}\par\vspace{1mm}\noindent\rule{\linewidth}{0.6pt}\par
Every key vocabulary word in this booklet, and the page where you wrote what it means.\par
\renewcommand{\arraystretch}{1.4}
\begin{longtable}{|L{6cm}|C{2cm}|}
\hline\thead{Word} & \thead{Page} \\ \hline \endhead
\glossrow{Aim}{voc:aim}
\glossrow{Anomaly}{voc:anomaly}
\glossrow{Average}{voc:average}
\glossrow{Boiling chip}{voc:boiling-chip}
\glossrow{Boiling point}{voc:boiling-point}
\glossrow{Cause and effect}{voc:cause-and-effect}
\glossrow{Condense}{voc:condense}
\glossrow{Crystal}{voc:crystal}
\glossrow{Crystallisation}{voc:crystallisation}
\glossrow{Decanting}{voc:decanting}
\glossrow{Disposal}{voc:disposal}
\glossrow{Dissolve}{voc:dissolve}
\glossrow{Distillate}{voc:distillate}
\glossrow{Distillation}{voc:distillation}
\glossrow{Ethical}{voc:ethical}
\glossrow{Evaporation}{voc:evaporation}
\glossrow{Filtrate}{voc:filtrate}
\glossrow{Filtration}{voc:filtration}
\glossrow{Flow chart}{voc:flow-chart}
\glossrow{Hazard}{voc:hazard}
\glossrow{Insoluble}{voc:insoluble}
\glossrow{Line of best fit}{voc:line-of-best-fit}
\glossrow{Magnetic separation}{voc:magnetic-separation}
\glossrow{Mass lost}{voc:mass-lost}
\glossrow{Method}{voc:method}
\glossrow{Mixture}{voc:mixture}
\glossrow{Pore}{voc:pore}
\glossrow{Property}{voc:property}
\glossrow{Rate}{voc:rate}
\glossrow{Rate of evaporation}{voc:rate-of-evaporation}
\glossrow{Repeatable}{voc:repeatable}
\glossrow{Residue}{voc:residue}
\glossrow{Risk}{voc:risk}
\glossrow{Saturated}{voc:saturated}
\glossrow{Seed crystal}{voc:seed-crystal}
\glossrow{Sieving}{voc:sieving}
\glossrow{Solubility}{voc:solubility}
\glossrow{Solubility curve}{voc:solubility-curve}
\glossrow{Soluble}{voc:soluble}
\glossrow{Solute}{voc:solute}
\glossrow{Solution}{voc:solution}
\glossrow{Solvent}{voc:solvent}
\glossrow{Surface area}{voc:surface-area}
\glossrow{Suspension}{voc:suspension}
\end{longtable}

\end{document}
